\documentclass[11pt]{amsart}

\usepackage[T1]{fontenc}
\usepackage{lmodern}
\usepackage[protrusion=true,expansion=false]{microtype}
\usepackage{amsmath,amssymb,amsthm,mathtools}
\usepackage[margin=2.5cm]{geometry}
\usepackage[colorlinks=true,linkcolor=blue,citecolor=blue,urlcolor=blue,
            hypertexnames=false]{hyperref}

\newtheorem{mainresult}{Theorem}
\renewcommand{\themainresult}{\Alph{mainresult}}
\newtheorem{maincorollary}[mainresult]{Corollary}
\newtheorem{theorem}{Theorem}[section]
\newtheorem{proposition}[theorem]{Proposition}
\newtheorem{lemma}[theorem]{Lemma}
\newtheorem{corollary}[theorem]{Corollary}
\theoremstyle{remark}
\newtheorem{remark}[theorem]{Remark}

\newcommand{\R}{\mathbb R}
\newcommand{\C}{\mathbb C}
\newcommand{\T}{\mathbb T}
\newcommand{\HH}{\mathcal H}
\newcommand{\Sc}{\mathcal S}
\newcommand{\supp}{\operatorname{supp}}
\newcommand{\Ran}{\operatorname{Ran}}
\newcommand{\ac}{\mathrm{ac}}
\newcommand{\pp}{\mathrm{pp}}
\newcommand{\Id}{\mathrm{Id}}
\newcommand{\D}{\mathcal D}
\newcommand{\1}{\mathbf 1}
\newcommand{\loc}{\mathrm{loc}}

\title{Phase retrieval for stationary Schr\"odinger evolutions}
\author{}
\date{}

\begin{document}
\raggedbottom

\begin{abstract}
We prove phase retrieval up to one constant unimodular factor for three
classes of stationary one-dimensional Schr\"odinger evolutions.  For real
potentials in the Faddeev class $L^1_1(\mathbb R)$, measurements on all of
spacetime determine arbitrary $L^2$ initial data.  If the potential is smooth
and compactly supported to the left of a point, measurements on the exterior
half-line to its right suffice.  Masuda's unique-continuation theorem gives a
third result for finite-energy solutions of a broad class of semibounded
potentials which may grow at infinity.  We also examine the attempt to extend
the Masuda argument to arbitrary $L^2$ data and isolate the missing current
trace.
\end{abstract}

\maketitle

\tableofcontents

\section*{Introduction and main results}

Let $V:\mathbb R\to\mathbb R$ and let $H=-\partial_x^2+V$ be a self-adjoint
Schr\"odinger operator.  We ask whether
\[
 |e^{-itH}f(x)|^2
\]
determines $f$ up to multiplication by one element of
$\mathbb T=\{z\in\mathbb C:|z|=1\}$.

Write
\[
 L^1_1(\mathbb R)=\left\{V:\mathbb R\to\mathbb C:
 \int_{\mathbb R}(1+|x|)|V(x)|\,dx<\infty\right\}.
\]
The first result allows arbitrary $L^2$ initial data, negative eigenvalues,
and a zero-energy resonance.

\setcounter{mainresult}{2}
\begin{mainresult}[Stationary potentials in the Faddeev class]
\label{thm:main}
Let $V\in L^1_1(\mathbb R)$ be real, let $H=-\partial_x^2+V$ be its
self-adjoint quadratic-form realization, and let $u_0,v_0\in L^2(\mathbb R)$.
If
\[
 |e^{-itH}u_0(x)|=|e^{-itH}v_0(x)|
 \quad\text{for almost every }(t,x)\in\mathbb R^2,
\]
then $v_0=\zeta u_0$ for some $\zeta\in\mathbb T$.
\end{mainresult}

For smooth compactly supported potentials, the observation set can be reduced
to one exterior half-line.

\begin{mainresult}[Observation on one exterior half-line]
\label{thm:half-line}
Let $V\in C_c^\infty(\mathbb R;\mathbb R)$, let $H=-\partial_x^2+V$, and
choose $R\in\mathbb R$ so that $\operatorname{supp}V\subset(-\infty,R]$.
If $u_0,v_0\in L^2(\mathbb R)$ and
\[
 |e^{-itH}u_0(x)|=|e^{-itH}v_0(x)|
 \quad\text{for almost every }(t,x)\in\mathbb R\times(R,\infty),
\]
then $v_0=\zeta u_0$ for some $\zeta\in\mathbb T$.
\end{mainresult}

The last theorem uses measurements on one nonempty open time interval and
allows semibounded potentials which need not decay.  Its additional $H^1$
hypothesis is part of the reduction from equal densities to an open zero set,
not an assumption in Masuda's unique-continuation theorem itself.

\begin{mainresult}[Masuda class]
\label{thm:phase-retrieval}
Let $V:\mathbb R\to\mathbb R$ satisfy the following conditions:
\begin{enumerate}
\item $V\in L^2_{\mathrm{loc}}(\mathbb R)$;
\item there is a null set $Q$ such that $\mathbb R\setminus Q$ is connected
and open and $V$ is locally essentially bounded on $\mathbb R\setminus Q$;
\item $-\partial_x^2+V$ on $C_c^\infty(\mathbb R)$ is essentially
self-adjoint;
\item its self-adjoint closure is bounded from below.
\end{enumerate}
Let $u,v\in C(\mathbb R;L^2(\mathbb R))$ be weak solutions of
$if_t+f_{xx}=Vf$.  Suppose that, on a nonempty open interval $I$,
\[
 u,v\in C(I;H^1(\mathbb R)),
 \qquad |u|=|v|\quad\text{almost everywhere on }I\times\mathbb R.
\]
Then $v=\zeta u$ on $\mathbb R\times\mathbb R$ for some
$\zeta\in\mathbb T$.
\end{mainresult}

Theorems~\ref{thm:main} and~\ref{thm:phase-retrieval} are incomparable.
The first has arbitrary $L^2$ data and an integrable decaying potential; the
second has finite-energy data but permits, for example, confining potentials.
Theorem~\ref{thm:half-line} strengthens the observation geometry on a smaller
potential class.

The proof of Theorem~\ref{thm:main} uses the one-dimensional scattering
representation in Proposition~\ref{prop:scattering}.  Theorem
\ref{thm:half-line} uses the same scattering data and a Cauchy-transform
argument.  Theorem~\ref{thm:phase-retrieval} uses Masuda's
unique-continuation theorem.  These external inputs are stated in the forms
needed below.

All nonsmooth identities are distributional unless another meaning is
specified.  We use
\[
 \widehat f(\xi)=(2\pi)^{-1/2}\int_{\mathbb R}e^{-ix\xi}f(x)\,dx
\]
and take Hilbert-space inner products to be linear in the first variable.
The value of a spectral representative on a null set is immaterial.

\part{Stationary decaying potentials}

\section{Stationary potentials in the Faddeev class}

Write
\[
 L^1_1(\R)=\left\{V:\R\to\C:
 \int_{\R}(1+|x|)|V(x)|\,dx<\infty\right\}.
\]
We prove Theorem~\ref{thm:main}.  Let $H=-\partial_x^2+V$ be the
self-adjoint quadratic-form realization, and choose jointly measurable
representatives of
\[
 u(t)=e^{-itH}u_0,
 \qquad
 v(t)=e^{-itH}v_0.
\]
The hypothesis is
\[
 |u(t,x)|=|v(t,x)|
 \quad\text{for almost every $(t,x)\in\R^2$}.
 \tag{1.1}\label{eq:modulus-hypothesis}
\]

Strong continuity of the unitary group and
\[
 \||a|^2-|b|^2\|_1
 \le\|a-b\|_2(\|a\|_2+\|b\|_2)
\]
show that $t\mapsto |e^{-itH}f|^2$ is continuous from $\R$ to $L^1(\R)$.
Thus \eqref{eq:modulus-hypothesis}, which initially holds for almost every
time, implies equality in $L^1_x$ for every $t$.  The proof uses scattering at both
infinities.  A zero-energy resonance is allowed.  The rank-one structure of
the spectral density matrix gives integrability on almost every fixed
energy-gap curve, while the product system in
Lemma~\ref{lem:product-system} uses no derivative of $V$.

\section{Scattering input}

We record the standard one-dimensional facts used in the proof.  They are
valid for real $V\in L^1_1(\R)$; see, for example,
\cite{DeiftTrubowitz,EgorovaKopylovaMarchenkoTeschl,TeschlBook}.

\begin{proposition}[Standard scattering representation]\label{prop:scattering}
Let $V\in L^1_1(\R)$ be real.  Then the following statements hold.
\begin{enumerate}
\item The singular continuous spectrum of $H$ is empty, the absolutely
continuous spectrum is $[0,\infty)$ with multiplicity two, and the negative
spectrum consists of finitely many simple eigenvalues
\[
 E_j=-\kappa_j^2,
 \qquad 0<\kappa_1<\cdots<\kappa_N.
\]
There is no eigenvalue in $[0,\infty)$.

\item There are real normalized eigenfunctions $\phi_j$ and physical
generalized eigenfunctions $e_\alpha(\cdot,k)$, $\alpha\in\{1,2\}$, $k>0$,
such that
\[
 \mathcal F:
 L^2(\R)\longrightarrow
 \C^N\oplus L^2(\R_+;\C^2)
\]
is unitary and diagonalizes $H$.  Thus, if
\[
 \mathcal F f=(c_1,\ldots,c_N;g_1,g_2),
\]
then
\begin{align}
 (e^{-itH}f)(x)
 ={}&\sum_{j=1}^N c_j e^{-itE_j}\phi_j(x) \notag\\
 &+\sum_{\alpha=1}^2\int_0^\infty
       e^{-itk^2}g_\alpha(k)e_\alpha(x,k)\,dk
 \tag{2.1}\label{eq:spectral-expansion}
\end{align}
in $L^2_x$.

\item For each $k>0$, the two functions
$e_1(\cdot,k),e_2(\cdot,k)$ form a basis of the solution space of
\[
 -y''+Vy=k^2y.
\]
On compact $x$-sets, uniformly for $k>0$,
\[
 |e_\alpha(x,k)|\lesssim 1,
 \qquad
 |\partial_xe_\alpha(x,k)|\lesssim 1+k.
 \tag{2.2}\label{eq:scattering-local-bounds}
\]
More precisely, for each $\sigma\in\{+,-\}$ there are coefficients
$c^{\sigma}_{\alpha,\varepsilon}(k)$, $\varepsilon\in\{+1,-1\}$, and
remainders $r^\sigma_\alpha(x,k)$ such that
\begin{equation}
 e_\alpha(x,k)
 =\sum_{\varepsilon=\pm1}
   c^{\sigma}_{\alpha,\varepsilon}(k)e^{i\varepsilon kx}
   +r^\sigma_\alpha(x,k),
 \qquad x\to\sigma\infty,
 \tag{2.3}\label{eq:uniform-scattering-expansion}
\end{equation}
and
\[
 \sup_{k>0}|c^{\sigma}_{\alpha,\varepsilon}(k)|<\infty,
 \qquad
 \lim_{x\to\sigma\infty}
 \sup_{k>0}|r^\sigma_\alpha(x,k)|=0.
\]

\item For every negative eigenfunction,
\[
 \phi_j(x)=a_j^+e^{-\kappa_jx}(1+o(1))
 \quad (x\to+\infty),
 \tag{2.4}\label{eq:bound-tail-plus}
\]
and there is an analogous expansion at $-\infty$, with nonzero leading
coefficients.
\end{enumerate}
\end{proposition}

For clarity, let $f_\pm(x,k)\sim e^{\pm ikx}$ as
$x\to\pm\infty$ be the Jost solutions.  We choose the physical incidence
eigenfunctions in the proposition to be
\[
 e_1(x,k)=(2\pi)^{-1/2}T(k)f_+(x,k),
 \qquad
 e_2(x,k)=(2\pi)^{-1/2}T(k)f_-(x,k).
\]
Thus the uniform assertion in part (3) concerns the physical incidence
eigenfunctions, not an unnormalized Jost solution on the opposite half-line.
The right Jost Volterra equation and the transformation-operator estimate give
\[
 \sup_{k>0}|f_+(x,k)e^{-ikx}-1|\longrightarrow0
 \qquad (x\to+\infty),
\]
and the analogous assertion holds for $f_-$ at $-\infty$.  Multiplication by
the transmission coefficient, followed by the scattering identity
\[
 T(k)f_+(x,k)=f_-(x,-k)+R_-(k)f_-(x,k),
\]
transfers the bounds to the left because $|T(k)|,|R_-(k)|\le1$.  The other
channel is analogous; here one also uses
$f_\pm(x,-k)=\overline{f_\pm(x,k)}$ for real $V$.  This proves
\eqref{eq:uniform-scattering-expansion}.  The same Volterra estimates, together
with the differential equation, give the local derivative bound in
\eqref{eq:scattering-local-bounds}.  This also shows directly why a zero
resonance causes no loss of uniformity for the physical eigenfunctions.

Only this proposition is imported from scattering theory.  Everything below is
an elementary consequence of it, one-dimensional ODE theory, and Fourier
analysis.

\section{The fixed-gap shell estimate}

Let
\[
 \delta=k^2-\ell^2,
 \qquad
 s=k^2+\ell^2,
 \qquad k,\ell>0.
\]
The inverse change of variables is
\[
 k=\sqrt{\frac{s+\delta}{2}},
 \qquad
 \ell=\sqrt{\frac{s-\delta}{2}},
 \qquad s>|\delta|,
\]
and
\[
 dk\,d\ell
 =\frac{d\delta\,ds}{4\sqrt{s^2-\delta^2}}.
 \tag{3.1}\label{eq:gap-sum-jacobian}
\]

The next estimate is the reason arbitrary $L^2$ initial states are admissible.

\begin{lemma}[Rank-one shell estimate]\label{lem:shell}
Let $f,g\in L^2(\R_+)$ and let $I\subset\R$ be a bounded interval.  Then
\[
 \iint_{\{k^2-\ell^2\in I\}}
 |f(k)g(\ell)|\,dk\,d\ell
 \le C_I\|f\|_2\|g\|_2.
 \tag{3.2}\label{eq:shell-schur}
\]
Consequently, for almost every $\delta\in\R$,
\[
 s\longmapsto
 \frac{
 |f(\sqrt{(s+\delta)/2})
 g(\sqrt{(s-\delta)/2})|
 }{4\sqrt{s^2-\delta^2}}
 \tag{3.3}\label{eq:shell-density}
\]
belongs to $L^1((|\delta|,\infty))$.  The resulting $L^1$ norm is locally
integrable as a function of $\delta$.
\end{lemma}

\begin{proof}
Set
\[
 K_I(k,\ell)=\mathbf 1_I(k^2-\ell^2),
 \qquad k,\ell>0.
\]
Choose $M>0$ with $I\subset[-M,M]$.  For fixed $k$, the condition
$K_I(k,\ell)\ne0$ implies
\[
 \max(0,k^2-M)\le \ell^2\le k^2+M.
\]
If $k^2\le M$, the length of this interval in the $\ell$ variable is at most
$\sqrt{k^2+M}\le\sqrt{2M}$.  If $k^2>M$, its length is
\begin{align*}
 \sqrt{k^2+M}-\sqrt{k^2-M}
 &=\frac{2M}{\sqrt{k^2+M}+\sqrt{k^2-M}}\\
 &\le \sqrt{2M}.
\end{align*}
Consequently
\[
 \sup_{k>0}\int_0^\infty K_I(k,\ell)\,d\ell\le\sqrt{2M}.
\]
Interchanging $k$ and $\ell$ gives the same estimate for the other Schur
integral.  Schur's test states that a nonnegative kernel $K$ satisfying
\[
 \sup_k\int K(k,\ell)\,d\ell\le A,
 \qquad
 \sup_\ell\int K(k,\ell)\,dk\le B
\]
defines an $L^2$ operator of norm at most $\sqrt{AB}$.  We therefore have
$\|K_I\|_{L^2\to L^2}\le\sqrt{2M}$ and
\begin{align*}
 \iint K_I(k,\ell)|f(k)g(\ell)|\,dk\,d\ell
 &=\int_0^\infty |f(k)|(K_I|g|)(k)\,dk\\
 &\le\|f\|_2\|K_I|g|\|_2
 \le\sqrt{2M}\|f\|_2\|g\|_2.
\end{align*}
This proves \eqref{eq:shell-schur}.

For $(k,\ell)\in(0,\infty)^2$, the map
$(k,\ell)\mapsto(\delta,s)=(k^2-\ell^2,k^2+\ell^2)$ is a $C^1$
diffeomorphism onto $\{(\delta,s):s>|\delta|\}$.  The change-of-variables
formula \eqref{eq:gap-sum-jacobian} gives, for every bounded interval $I$,
\begin{align*}
 &\iint_{\{k^2-\ell^2\in I\}}|f(k)g(\ell)|\,dk\,d\ell\\
 &\quad=\int_I\int_{|\delta|}^\infty
 \frac{|f(\sqrt{(s+\delta)/2})g(\sqrt{(s-\delta)/2})|}
 {4\sqrt{s^2-\delta^2}}\,ds\,d\delta.
\end{align*}
The left side is finite by \eqref{eq:shell-schur}.  Tonelli's theorem now
shows that the inner integral is finite for almost every $\delta\in I$, and
that these inner integrals form an $L^1(I)$ function.  A countable exhaustion
of $\R$ by bounded intervals proves both remaining assertions.
\end{proof}

\begin{remark}\label{rem:rank-one-essential}
It would not be enough merely to observe that the spectral density-matrix
kernel belongs to $L^2(\R_+^2)$.  A general $L^2$ kernel need not have an
integrable restriction to almost every hyperbola $k^2-\ell^2=\delta$.
Lemma~\ref{lem:shell} uses the finite sum of rank-one terms in an essential
way.
\end{remark}

We will also use the following consequence of the one-dimensional change of
variables.  Suppose
$g\in L^2(\R_+)$ and $E_j=-\kappa_j^2$.  A bound--continuous term has temporal
frequency $k^2+\kappa_j^2$.  Its density in that frequency variable is locally
integrable, because under $\delta=k^2+\kappa_j^2$,
\[
 \frac{|g(\sqrt{\delta-\kappa_j^2})|}
      {2\sqrt{\delta-\kappa_j^2}}\,d\delta
 =|g(k)|\,dk,
\]
and $g$ is integrable on bounded $k$-intervals.  Bound--bound terms, by
contrast, are point masses in the temporal-frequency variable.

\section{A derivative-free product system}

The following elementary system is the central structural observation.

\begin{lemma}[Product system]\label{lem:product-system}
Let $V\in L^1_{\mathrm{loc}}(\R)$ and suppose
\[
 -y''+Vy=Ey,
 \qquad
 -z''+Vz=Fz
\]
distributionally.  Put
\[
 \delta=E-F,
 \qquad
 s=E+F,
\]
and define
\[
 A=yz,
 \quad
 B=y'z+yz',
 \quad
 C=y'z-yz',
 \quad
 D=y'z'.
 \tag{4.1}\label{eq:product-state-definition}
\]
Then $A,B,C,D$ are locally absolutely continuous and satisfy
\begin{equation}
 \boxed{
 \begin{aligned}
 A'&=B,\\
 B'&=(2V-s)A+2D,\\
 C'&=-\delta A,\\
 D'&=\left(V-\frac{s}{2}\right)B+\frac{\delta}{2}C.
 \end{aligned}}
 \tag{4.2}\label{eq:product-system}
\end{equation}
\end{lemma}

\begin{proof}
Since $V\in L^1_{\mathrm{loc}}$, every distributional solution belongs to
$W^{2,1}_{\mathrm{loc}}$.  Thus $y,z$ and their first derivatives are locally
absolutely continuous, and the ordinary product rule is valid almost
everywhere.  Since
\[
 y''=(V-E)y,
 \qquad z''=(V-F)z,
\]
we obtain
\[
 A'=y'z+yz'=B.
\]
Differentiating $B$ gives
\begin{align*}
 B'
 &=y''z+2y'z'+yz''\\
 &=(V-E)yz+2y'z'+(V-F)yz\\
 &=(2V-E-F)A+2D=(2V-s)A+2D.
\end{align*}
Similarly,
\begin{align*}
 C'
 &=y''z-yz''\\
 &=(V-E)yz-(V-F)yz=(F-E)A=-\delta A.
\end{align*}
Finally,
\begin{align*}
 D'
 &=y''z'+y'z''=(V-E)yz'+(V-F)y'z.
\end{align*}
The definitions of $B$ and $C$ imply
\[
 yz'=\frac{B-C}{2},
 \qquad y'z=\frac{B+C}{2}.
\]
Since $E=(s+\delta)/2$ and $F=(s-\delta)/2$, substitution gives
\begin{align*}
 D'
 &=\left(V-\frac{s+\delta}{2}\right)\frac{B-C}{2}
   +\left(V-\frac{s-\delta}{2}\right)\frac{B+C}{2}\\
 &=\left(V-\frac{s}{2}\right)B+\frac{\delta}{2}C.
\end{align*}
Every right-hand side is locally integrable.  Hence all four components are
locally absolutely continuous and the system holds everywhere in its
integral form.
\end{proof}

For fixed $E$ let $\Sc_E$ denote the solution space of
$-y''+Vy=Ey$.  Fix $x_0\in\R$.  Existence and uniqueness for the integral
system associated with the Schr\"odinger equation show that
\[
 y\longmapsto (y(x_0),y'(x_0))
\]
is an isomorphism from $\Sc_E$ onto $\C^2$.  In particular,
$\dim\Sc_E=2$.  The tensor product $\Sc_E\otimes\Sc_F$ is therefore
four-dimensional.  At $x_0$, a simple tensor $y\otimes z$ has coordinates
\[
 yz,\qquad y'z,\qquad yz',\qquad y'z'.
\]
The product-state coordinates are obtained from these by the invertible
linear change
\[
 A=yz,
 \quad B=y'z+yz',
 \quad C=y'z-yz',
 \quad D=y'z',
\]
whose inverse is
\[
 yz=A,
 \qquad y'z=\frac{B+C}{2},
 \qquad yz'=\frac{B-C}{2},
 \qquad y'z'=D.
\]
The first-order system \eqref{eq:product-system} has unique solutions for
arbitrary initial data in $\C^4$.  It follows that the map
\[
 \Sc_E\otimes\Sc_F
 \longrightarrow
 \{\text{solutions of \eqref{eq:product-system}}\}
 \tag{4.3}\label{eq:tensor-identification}
\]
induced by \eqref{eq:product-state-definition} is an isomorphism.

\begin{lemma}[Independence of unequal-energy products]
\label{lem:product-independence}
Suppose $E\ne F$.  If $y_1,y_2$ form a basis of $\Sc_E$ and $z_1,z_2$ form a
basis of $\Sc_F$, then the four functions
\[
 y_i z_j,
 \qquad 1\le i,j\le2,
\]
are linearly independent.
\end{lemma}

\begin{proof}
Suppose
\[
 \sum_{i,j=1}^2c_{ij}y_i(x)z_j(x)=0
 \quad\text{for every }x.
\]
The tensor $T=\sum_{i,j}c_{ij}y_i\otimes z_j$ determines, through
\eqref{eq:tensor-identification}, a solution $(A,B,C,D)$ of
\eqref{eq:product-system}.  The displayed relation says $A=0$.  Therefore
$B=A'=0$.  The second equation of the system becomes $0=2D$, so $D=0$.
The fourth equation then becomes $0=\delta C/2$.  Since
$\delta=E-F\ne0$, it follows that $C=0$.  Thus the product state is the zero
solution.  The map \eqref{eq:tensor-identification} is injective, so $T=0$.
Because $\{y_i\otimes z_j:1\le i,j\le2\}$ is a basis of
$\Sc_E\otimes\Sc_F$, every coefficient $c_{ij}$ vanishes.
\end{proof}

We also record the elementary negative-energy ODE fact used below.

\begin{lemma}[Negative-energy dichotomy]\label{lem:negative-dichotomy}
Let $V\in L^1(\R)$ be real and let $G=-\kappa^2<0$.  At $+\infty$ the
solution space $\Sc_G$ has a basis $y^d_{G,+},y^g_{G,+}$ satisfying
\begin{align*}
 y^d_{G,+}(x)&=e^{-\kappa x}(1+o(1)),
 & (y^d_{G,+})'(x)&=-\kappa e^{-\kappa x}(1+o(1)),\\
 y^g_{G,+}(x)&=e^{\kappa x}(1+o(1)),
 & (y^g_{G,+})'(x)&=\kappa e^{\kappa x}(1+o(1)).
\end{align*}
At $-\infty$ there is a basis $y^d_{G,-},y^g_{G,-}$ with
\begin{align*}
 y^d_{G,-}(x)&=e^{\kappa x}(1+o(1)),
 & (y^d_{G,-})'(x)&=\kappa e^{\kappa x}(1+o(1)),\\
 y^g_{G,-}(x)&=e^{-\kappa x}(1+o(1)),
 & (y^g_{G,-})'(x)&=-\kappa e^{-\kappa x}(1+o(1)).
\end{align*}
In particular, if the lines
\[
 L^d_{G,+}:=\operatorname{span}\{y^d_{G,+}\},
 \qquad
 L^d_{G,-}:=\operatorname{span}\{y^d_{G,-}\}
\]
coincide, then $G$ is an $L^2$ eigenvalue of $H$.
\end{lemma}

\begin{proof}
The equation at energy $G=-\kappa^2$ is
\[
 y''-\kappa^2y=Vy.
\]
On a right half-line, the solution which is asymptotic to $e^{-\kappa x}$ is
obtained from the Volterra equation
\[
 y(x)=e^{-\kappa x}
 +\frac1\kappa\int_x^\infty
   \sinh(\kappa(t-x))V(t)y(t)\,dt.
 \tag{4.4}\label{eq:negative-volterra}
\]
To see why only $V\in L^1$ is needed, write $m(x)=e^{\kappa x}y(x)$.  Then
\eqref{eq:negative-volterra} becomes
\[
 m(x)=1+\frac1{2\kappa}\int_x^\infty
 \bigl(1-e^{-2\kappa(t-x)}\bigr)V(t)m(t)\,dt.
\]
The kernel in parentheses has absolute value at most $2$.  Successive
approximation on a sufficiently far right half-line, followed by continuation
to the left, gives a unique bounded solution with $m(x)\to1$.  Differentiating
the Volterra equation gives $m'(x)\to0$.  Therefore
\[
 y^d_{G,+}(x)=e^{-\kappa x}(1+o(1)),
 \qquad
 (y^d_{G,+})'(x)=-\kappa e^{-\kappa x}(1+o(1)).
\]
In particular, $y^d_{G,+}$ has no zeros for all sufficiently large $x$.
Fix such an $x_0$.  Reduction of order gives a second solution
\[
 y^g_{G,+}(x)
 =2\kappa y^d_{G,+}(x)
   \int_{x_0}^{x}\frac{dt}{(y^d_{G,+}(t))^2}.
 \tag{4.5}\label{eq:negative-growing-solution}
\]
Since $(y^d_{G,+}(t))^{-2}=e^{2\kappa t}(1+o(1))$, integration yields
\[
 \int_{x_0}^{x}\frac{dt}{(y^d_{G,+}(t))^2}
 =\frac{e^{2\kappa x}}{2\kappa}(1+o(1)).
\]
Substitution in \eqref{eq:negative-growing-solution}, followed by
differentiation, gives the asserted growing asymptotics.  The Wronskian of
the two solutions is $2\kappa\ne0$, so they form a basis.  ODE uniqueness
extends both solutions to all of $\R$.

Replacing $x$ by $-x$ and $V(x)$ by $V(-x)$ gives the corresponding decaying
and growing basis at $-\infty$.

Suppose now that $L^d_{G,+}=L^d_{G,-}$.  A nonzero solution $y$ spanning this
common line satisfies, for some constants $C,R>0$,
\[
 |y(x)|+|y'(x)|\le Ce^{-\kappa|x|}
 \quad (|x|\ge R).
\]
Hence $y,y'\in L^2(\R)$, so $y\in H^1(\R)$.  The distributional equation
$-y''+Vy=Gy$ is the weak eigenvalue equation for the quadratic-form
realization of $H$.  Since its right side belongs to $L^2$, the representation
theorem for closed semibounded forms places $y$ in the operator domain and
gives $Hy=Gy$.  Thus $G$ is an $L^2$ eigenvalue.
\end{proof}

\section{Fixed-gap injectivity}

We now prove the analytic core of the argument.  Fix $\delta\ne0$.  For
$s>|\delta|$, put
\[
 E(s)=\frac{s+\delta}{2},
 \qquad
 F(s)=\frac{s-\delta}{2},
\]
and let
\[
 X_{\alpha\beta}(x,s)
 =\bigl(A_{\alpha\beta},B_{\alpha\beta},
        C_{\alpha\beta},D_{\alpha\beta}\bigr)(x,s)
\]
be the product state associated with
\[
 y=e_\alpha(\cdot,\sqrt{E(s)}),
 \qquad
 z=\overline{e_\beta(\cdot,\sqrt{F(s)})}.
\]
Since $V$ is real, the second factor is again a solution of the Schr\"odinger
equation at energy $F(s)$.

There may also be bound--continuous products with the same energy difference
$\delta$.  They form the following finite family.  If $\delta>\kappa_j^2$, set
\[
 k_j=\sqrt{\delta-\kappa_j^2},
 \qquad s_j=\delta-2\kappa_j^2,
\]
and take the two product states associated with
$e_\alpha(\cdot,k_j)\phi_j$, $\alpha=1,2$.  Their energies are
$k_j^2$ and $-\kappa_j^2$, so their difference is $\delta$ and their sum is
$s_j$.  If $\delta<-\kappa_j^2$, set
\[
 k_j=\sqrt{-\delta-\kappa_j^2},
 \qquad s_j=-\delta-2\kappa_j^2,
\]
and take the states associated with
$\phi_j\overline{e_\alpha(\cdot,k_j)}$.  These have energy difference
$\delta$ as well.  Enumerate all compatible states as $X_\rho(x)$, denote
their energy sums by $s_\rho$, and denote their scalar coefficients by
$a_\rho$.  If several contributions produce the same state, combine their
coefficients first.

\begin{proposition}[Fixed-gap injectivity]\label{prop:fixed-gap}
Suppose $q_{\alpha\beta}\in L^1((|\delta|,\infty))$ and
\begin{align}
 W(x)
 :={}&\sum_{\alpha,\beta=1}^2
 \int_{|\delta|}^\infty
 q_{\alpha\beta}(s)A_{\alpha\beta}(x,s)\,ds
 +\sum_\rho a_\rho A_\rho(x)
 =0
 \tag{5.1}\label{eq:fixed-gap-relation}
\end{align}
for every $x\in\R$.  Then every $q_{\alpha\beta}$ vanishes almost everywhere
and every $a_\rho$ vanishes.
\end{proposition}

\begin{proof}
The proof has four steps.

\smallskip
\noindent
\emph{Step 1: the energy-sum resolvent.}
The continuum sums satisfy $s>|\delta|$.  There are only finitely many
bound--continuous sums $s_\rho$.  Hence the set of all sums occurring in
\eqref{eq:fixed-gap-relation} is bounded from below.  Choose a real number
$\lambda$ strictly below this set.  In particular, every denominator below is
nonzero and has one sign.  Define
\begin{align}
 Z_\lambda(x)
 :={}&\sum_{\alpha,\beta=1}^2
 \int_{|\delta|}^\infty
 \frac{q_{\alpha\beta}(s)}{s-\lambda}
 X_{\alpha\beta}(x,s)\,ds
 +\sum_\rho\frac{a_\rho}{s_\rho-\lambda}X_\rho(x).
 \tag{5.2}\label{eq:resolvent-state}
\end{align}
Write its components as
$Z_\lambda=(A_\lambda,B_\lambda,C_\lambda,D_\lambda)$.

Let $K\subset\R$ be compact.  The energies $E(s)$ and $F(s)$ are positive and
at most $s+|\delta|$.  The local scattering estimates
\eqref{eq:scattering-local-bounds} therefore give, uniformly for $x\in K$ and
$s>|\delta|$,
\[
 |A_{\alpha\beta}|\\lesssim_K1,
 \qquad
 |B_{\alpha\beta}|+|C_{\alpha\beta}|
 \lesssim_K1+\sqrt{s},
 \qquad
 |D_{\alpha\beta}|\lesssim_K1+s.
 \tag{5.3}\label{eq:product-state-bounds}
\]
For example, $D$ is the product of one derivative at energy $E(s)$ and one at
energy $F(s)$, so
\[
 |D|\lesssim_K(1+\sqrt{E(s)})(1+\sqrt{F(s)})\lesssim_K1+s.
\]
Because $\lambda$ is fixed below the spectral sums,
\[
 \frac1{|s-\lambda|}\lesssim_\lambda\frac1{1+s}.
\]
After division by $s-\lambda$, the bounds for the four components are
respectively $O(1)$, $O((1+s)^{-1/2})$, $O((1+s)^{-1/2})$, and $O(1)$.
Each is bounded by a constant times $|q_{\alpha\beta}(s)|$.  Since every
$q_{\alpha\beta}$ is integrable, all components in
\eqref{eq:resolvent-state} converge absolutely and uniformly on $K$.  The
finite atomic sum causes no convergence issue.

Direct differentiation under the untruncated integral would require more
than $q\in L^1$, because derivatives of $D$ grow with $s$.  We avoid this.
Choose $\chi_R\in C_c^\infty((|\delta|,\infty))$ with
$0\le\chi_R\le1$ and $\chi_R(s)\to1$ pointwise.  Insert $\chi_R$ in every
continuum integral in \eqref{eq:resolvent-state}, retain all atomic terms, and
call the resulting state $Z_{\lambda,R}$.  Let $W_R$ be the same truncation of
\eqref{eq:fixed-gap-relation}.  Every truncated integral may be
differentiated distributionally by pairing with a compactly supported test
function in $x$ and using Fubini's theorem.  Applying
\eqref{eq:product-system} and the identity
\[
 \frac{s}{s-\lambda}=1+\frac{\lambda}{s-\lambda}
\]
gives
\begin{align}
 A_{\lambda,R}'
 &=B_{\lambda,R}, \tag{5.4}\label{eq:cutoff-system-1}\\
 B_{\lambda,R}'
 &=(2V-\lambda)A_{\lambda,R}+2D_{\lambda,R}-W_R,
 \tag{5.5}\label{eq:cutoff-system-2}\\
 C_{\lambda,R}'
 &=-\delta A_{\lambda,R}, \tag{5.6}\label{eq:cutoff-system-3}\\
 D_{\lambda,R}'
 &=\left(V-\frac{\lambda}{2}\right)B_{\lambda,R}
   +\frac{\delta}{2}C_{\lambda,R}-\frac12W_R'.
 \tag{5.7}\label{eq:cutoff-system-4}
\end{align}
The uniform local bound for $A$ and dominated convergence give
$W_R\to W=0$ locally uniformly.  Therefore, for every
$\varphi\in C_c^\infty(\R)$,
\[
 \langle W_R',\varphi\rangle
 =-\int_\R W_R(x)\varphi'(x)\,dx\longrightarrow0.
\]
Thus $W_R'\to0$ in distributions.  The component bounds established above
also give local uniform convergence $Z_{\lambda,R}\to Z_\lambda$.  Products
such as $V B_{\lambda,R}$ converge in $L^1_{\mathrm{loc}}$, because
$V\in L^1_{\mathrm{loc}}$ and $B_{\lambda,R}$ converges locally uniformly.
We may therefore pass to the limit in each of the four distributional
identities and obtain
\begin{equation}
 \boxed{
 \begin{aligned}
 A_\lambda'&=B_\lambda,\\
 B_\lambda'&=(2V-\lambda)A_\lambda+2D_\lambda,\\
 C_\lambda'&=-\delta A_\lambda,\\
 D_\lambda'&=\left(V-\frac{\lambda}{2}\right)B_\lambda
              +\frac{\delta}{2}C_\lambda.
 \end{aligned}}
 \tag{5.8}\label{eq:lambda-product-system}
\end{equation}
Initially these identities hold distributionally, while all four components
are continuous.  The first and third equations already show that $A_\lambda$
and $C_\lambda$ are locally absolutely continuous.  The right side of the
fourth equation belongs to $L^1_{\mathrm{loc}}$: $V B_\lambda$ is locally
integrable and the other terms are continuous.  Hence $D_\lambda$ is locally
absolutely continuous.  The right side of the second equation is then locally
integrable, so $B_\lambda$ is locally absolutely continuous as well.  Thus
$Z_\lambda$ is an ordinary solution of the product system with auxiliary
energies
\[
 E_\lambda=\frac{\lambda+\delta}{2},
 \qquad
 F_\lambda=\frac{\lambda-\delta}{2}.
 \tag{5.9}\label{eq:auxiliary-energies}
\]

\smallskip
\noindent
\emph{Step 2: decay of the scalar component.}
We claim that
\[
 A_\lambda(x)\longrightarrow0
 \qquad (x\to\pm\infty).
 \tag{5.10}\label{eq:scalar-decay}
\]
Every atomic term is a product of a fixed positive-energy scattering solution
and a negative-energy eigenfunction.  The scattering factor is bounded at
both ends, while the eigenfunction decays exponentially there.  Hence every
atomic contribution to $A_\lambda(x)$ tends to zero as $x\to\pm\infty$.

Consider a continuum product at one of the two ends.  Write both scattering
solutions as the sum of the two plane waves in
\eqref{eq:uniform-scattering-expansion} and a remainder.  Multiplication gives
four plane-wave terms, with phases
\[
 x\bigl(\varepsilon k(s)+\eta\ell(s)\bigr),
 \qquad \varepsilon,\eta\in\{-1,1\},
 \tag{5.11}\label{eq:product-phases}
\]
and terms containing at least one remainder.  The plane-wave coefficients are
uniformly bounded.  Therefore their amplitudes after insertion into
$A_\lambda$ are bounded by
\[
 C\frac{|q_{\alpha\beta}(s)|}{|s-\lambda|},
\]
which belongs to $L^1((|\delta|,\infty))$.  The scattering remainders tend to
zero uniformly in the energy parameter.  Their products are bounded by the
same integrable majorant, so dominated convergence makes their integrals tend
to zero.

It remains to treat each plane-wave integral.  Define
\[
 p_+(s)=k(s)+\ell(s),
 \qquad
 p_-(s)=k(s)-\ell(s).
\]
Since $k^2-\ell^2=\delta$,
\[
 p_-(s)=\frac{\delta}{p_+(s)}.
\]
The function $p_+$ is strictly increasing: both $k(s)$ and $\ell(s)$ are
strictly increasing.  Since $\delta\ne0$, $p_-$ is strictly monotone as well.
Their derivatives do not vanish on the open interval
$(|\delta|,\infty)$.  Thus each is a $C^1$ diffeomorphism onto its image.
If $f\in L^1(ds)$, the substitution $p=p_\pm(s)$ gives
\[
 \int_{|\delta|}^\infty f(s)e^{ixp_\pm(s)}\,ds
 =\int_{p_\pm((|\delta|,\infty))}
 f(s(p))\left|\frac{ds}{dp}\right|e^{ixp}\,dp.
\]
The new amplitude is integrable because its $L^1(dp)$ norm equals
$\|f\|_{L^1(ds)}$.  The Riemann--Lebesgue lemma makes the last integral tend
to zero as $|x|\to\infty$.  The four phases in
\eqref{eq:product-phases} are $\pm p_+$ and $\pm p_-$.  Summing them proves
\eqref{eq:scalar-decay} at both ends.

\smallskip
\noindent
\emph{Step 3: uniqueness at a negative auxiliary sum.}
Choose $\lambda$ sufficiently negative that
\[
 \lambda<-|\delta|,
 \qquad
 \frac{\lambda-|\delta|}{2}<\inf\sigma(H).
 \tag{5.12}\label{eq:negative-lambda-choice}
\]
Then $E_\lambda,F_\lambda<0$, and the lower of these two energies lies below
the spectrum of $H$.

Assume first that $\delta>0$.  Write
\[
 E_\lambda=-\alpha^2,
 \qquad
 F_\lambda=-\beta^2,
 \qquad 0<\alpha<\beta.
\]
By Lemma~\ref{lem:negative-dichotomy}, choose decaying and growing bases at
$+\infty$ for both energy spaces.  In these bases the tensor represented by
$Z_\lambda$ has a unique expansion
\begin{align*}
 T={}&c_{dd}\,y^d_{E_\lambda,+}\otimes y^d_{F_\lambda,+}
 +c_{gd}\,y^g_{E_\lambda,+}\otimes y^d_{F_\lambda,+}\\
 &+c_{dg}\,y^d_{E_\lambda,+}\otimes y^g_{F_\lambda,+}
 +c_{gg}\,y^g_{E_\lambda,+}\otimes y^g_{F_\lambda,+}.
\end{align*}
Its scalar component $A_\lambda$ is the same linear combination of the four
ordinary products.  Their asymptotic rates are, in the displayed order,
\[
 e^{-(\alpha+\beta)x},
 \quad e^{-(\beta-\alpha)x},
 \quad e^{(\beta-\alpha)x},
 \quad e^{(\alpha+\beta)x}.
\]
Divide $A_\lambda(x)$ by $e^{(\alpha+\beta)x}$.  Every term except the last
tends to zero, while the last tends to $c_{gg}$.  Since
$A_\lambda(x)\to0$, we obtain $c_{gg}=0$.  Now divide the remaining expansion
by $e^{(\beta-\alpha)x}$.  The first two terms tend to zero and the third
tends to $c_{dg}$, so $c_{dg}=0$.  The coefficients $c_{dd}$ and $c_{gd}$ are
not forced to vanish at this end because their products decay.  We have shown
that
\[
 T\in\Sc_{E_\lambda}\otimes L^d_{F_\lambda,+}.
\]
where $L^d_{F_\lambda,+}$ is the one-dimensional space of solutions at energy
$F_\lambda$ which decay at $+\infty$.  Repeating the argument at $-\infty$
places the same tensor in
\[
 \Sc_{E_\lambda}\otimes L^d_{F_\lambda,-}.
\]
The two lines $L^d_{F_\lambda,+}$ and $L^d_{F_\lambda,-}$ are distinct by
Lemma~\ref{lem:negative-dichotomy}; otherwise $F_\lambda$ would be an $L^2$
eigenvalue, contrary to $F_\lambda<\inf\sigma(H)$.  Distinct one-dimensional subspaces of the
two-dimensional space $\Sc_{F_\lambda}$ have zero intersection.  To verify
the tensor statement, choose a basis $z_+,z_-$ of $\Sc_{F_\lambda}$ with
$z_+$ spanning $L^d_{F_\lambda,+}$ and $z_-$ spanning
$L^d_{F_\lambda,-}$.  A tensor in the first product space has the form
$Y\otimes z_+$, while one in the second has the form $\widetilde Y\otimes
z_-$.  Uniqueness of coordinates in the basis $z_+,z_-$ forces both tensors
to vanish if they are equal.  Hence the two tensor-product subspaces have zero
intersection, so $T=0$ and therefore $Z_\lambda=0$.  If
$\delta<0$, the same argument applies after interchanging $E_\lambda$ and
$F_\lambda$.

We have proved
\[
 Z_\lambda=0
 \quad\text{for every sufficiently negative real $\lambda$}.
 \tag{5.13}\label{eq:negative-resolvent-zero}
\]

\smallskip
\noindent
\emph{Step 4: Cauchy uniqueness.}
For fixed $x$, define the finite complex measure
\begin{align}
 d\mu_x(s)
 :={}&\sum_{\alpha,\beta=1}^2
 q_{\alpha\beta}(s)A_{\alpha\beta}(x,s)\,ds
 +\sum_\rho a_\rho A_\rho(x)\delta_{s_\rho}.
 \tag{5.14}\label{eq:cauchy-measure}
\end{align}
Choose $m\in\R$ strictly below every spectral sum occurring in
\eqref{eq:fixed-gap-relation}; in particular
$\supp\mu_x\subset[m,\infty)$.  Set
\[
 F_x(z):=\int_{[m,\infty)}\frac{d\mu_x(s)}{s-z},
 \qquad z\in\C\setminus[m,\infty).
\]
For every compact subset of the slit plane, the denominator is uniformly
separated from zero on the support of the measure.  Differentiation under the
integral is therefore justified, so $F_x$ is holomorphic on the connected
domain $\C\setminus[m,\infty)$.  For real $\lambda<m$, its defining integral
is exactly the first component $A_\lambda(x)$ of
\eqref{eq:resolvent-state}.  By \eqref{eq:negative-resolvent-zero}, it
vanishes on a nonempty real interval.  The holomorphic identity theorem gives
$F_x\equiv0$ on the slit plane.

We recall why this forces the measure to vanish.  For $a<b$ at which
$\mu_x$ has no atoms, the Stieltjes inversion formula gives
\[
 \mu_x((a,b))
 =\lim_{\varepsilon\downarrow0}\frac1{2\pi i}
 \int_a^b\bigl(F_x(t+i\varepsilon)-F_x(t-i\varepsilon)\bigr)\,dt.
\]
The right side is zero.  Applying this separately to the real and imaginary
parts and then using regularity of finite Borel measures gives $\mu_x=0$.

For a fixed $x$, the equality $\mu_x=0$ says that its absolutely continuous
density vanishes for almost every $s$, and that each of its atomic masses is
zero.  The exceptional set in $s$ may depend on $x$.  Choose a countable dense
set $D\subset\R$.  The union of the exceptional sets over $x\in D$ is still
null.  For every $s$ outside this union,
\[
 \sum_{\alpha,\beta=1}^2
 q_{\alpha\beta}(s)A_{\alpha\beta}(x,s)=0
 \quad (x\in D).
\]
For fixed $s$, each product is continuous in $x$.  Density of $D$ extends the
identity to every $x\in\R$.
Lemma~\ref{lem:product-independence} gives
$q_{\alpha\beta}(s)=0$ for every $\alpha,\beta$.  At an atom, distinct bound
energies give distinct values of $s$ for fixed $\delta$.  More explicitly, if
$\delta>0$ and the second energy is $E_j=-\kappa_j^2$, then
\[
 E=\delta-\kappa_j^2>0,
 \qquad
 s=\delta-2\kappa_j^2,
 \qquad |s|<\delta.
\]
Thus this atom lies strictly below the continuum threshold $s=\delta$, and
different $j$ give different atoms.  A zero residue has the form
\[
 \phi_j(x)\sum_{\alpha=1}^2b_\alpha e_\alpha(x,k)=0.
\]
The eigenfunction $\phi_j$ is nonzero on some interval, and the two continuum
solutions are independent there, so $b_1=b_2=0$.  The case $\delta<0$ is
symmetric.  Thus all atomic coefficients vanish as well.
\end{proof}

\section{Application to the modulus identity}

Let
\[
 \mathcal Fu_0=(c_1,\ldots,c_N;g_1,g_2),
 \qquad
 \mathcal Fv_0=(d_1,\ldots,d_N;h_1,h_2).
\]
Choose measurable representatives of the four $L^2$ functions; their values
on null sets will play no role.
Define
\[
 Q_{\alpha\beta}(k,\ell)
 =g_\alpha(k)\overline{g_\beta(\ell)}
  -h_\alpha(k)\overline{h_\beta(\ell)}.
 \tag{6.1}\label{eq:ac-kernel-difference}
\]
This is a finite sum of rank-one $L^2$ kernels.  Formally, the
continuous--continuous part of the density difference is
\begin{align*}
 \sum_{\alpha,\beta=1}^2\int_0^\infty\int_0^\infty
 &e^{-it(k^2-\ell^2)}Q_{\alpha\beta}(k,\ell)\\
 &\times e_\alpha(x,k)\overline{e_\beta(x,\ell)}\,dk\,d\ell.
 \tag{6.2a}\label{eq:formal-cc-density}
\end{align*}
The temporal frequency is the energy difference
$\delta=k^2-\ell^2$.  The second variable
$s=k^2+\ell^2$ parametrizes the corresponding fixed-gap curve.  The purpose
of the next approximation is to make this formal expansion valid for
$L^2$ spectral coefficients.

The equality of moduli says that
\[
 \mathcal D(t,x):=|u(t,x)|^2-|v(t,x)|^2=0
 \tag{6.2}\label{eq:density-difference}
\]
as a distribution in $(t,x)$.  To justify the spectral expansion, choose
$g_{\alpha,n},h_{\alpha,n}\in C_c^\infty((0,\infty))$ converging to
$g_\alpha,h_\alpha$ in $L^2$, while leaving the bound-state coefficients
fixed.  Let $u_n,v_n$ be the corresponding evolutions.  Unitarity gives
\[
 \sup_{t\in\R}\|u_n(t)-u(t)\|_2=\|u_{0,n}-u_0\|_2\longrightarrow0.
\]
For any $a,b\in L^2$, the identity
$|a|^2-|b|^2=(a-b)\overline a+b(\overline a-\overline b)$ and
Cauchy--Schwarz give
\[
 \||a|^2-|b|^2\|_1\le\|a-b\|_2(\|a\|_2+\|b\|_2).
\]
Consequently, uniformly in $t$,
\[
 \bigl\||u_n(t)|^2-|u(t)|^2\bigr\|_{L^1_x}
 \le
 \|u_n(t)-u(t)\|_2\bigl(\|u_n(t)\|_2+\|u(t)\|_2\bigr)
 \longrightarrow0,
 \tag{6.3}\label{eq:density-approximation}
\]
and the same holds for $v_n$.

Fix a spatial test function $\varphi\in C_c^\infty(\R)$ and a temporal
frequency cutoff $\chi\in C_c^\infty((-M,M))$.  On $\supp\varphi$, the
physical eigenfunctions are uniformly bounded.  Thus the difference between
the $(\alpha,\beta)$ continuous--continuous forms for $u_n$ and $u$ is
bounded, up to a constant depending on $\varphi$, by
\begin{align*}
 \iint_{|k^2-\ell^2|<M}
 &\bigl|g_{\alpha,n}(k)\overline{g_{\beta,n}(\ell)}
       -g_\alpha(k)\overline{g_\beta(\ell)}\bigr|\,dk\,d\ell\\
 \le{}&\iint_{|k^2-\ell^2|<M}
 |g_{\alpha,n}(k)-g_\alpha(k)|\,|g_{\beta,n}(\ell)|\,dk\,d\ell\\
 &+\iint_{|k^2-\ell^2|<M}
 |g_\alpha(k)|\,|g_{\beta,n}(\ell)-g_\beta(\ell)|\,dk\,d\ell.
\end{align*}
Lemma~\ref{lem:shell} bounds this by a constant times
\[
 \sqrt{2M}\left(
 \|g_{\alpha,n}-g_\alpha\|_2\|g_{\beta,n}\|_2
 +\|g_\alpha\|_2\|g_{\beta,n}-g_\beta\|_2
 \right),
 \tag{6.4}\label{eq:double-form-approximation}
\]
which tends to zero.  The same estimate applies to $h$.

Let $\check\chi$ be the inverse Fourier transform of $\chi$.  Since
$\check\chi$ is Schwartz, it belongs to $L^1(\R_t)$.  Equation
\eqref{eq:density-approximation} therefore permits passage to the limit in
\[
 \iint \check\chi(t)\varphi(x)
 \bigl(|u_n(t,x)|^2-|v_n(t,x)|^2\bigr)\,dx\,dt.
\]
For the smooth compactly supported spectral data, Fubini's theorem gives the
spectral expansion directly.  Estimate \eqref{eq:double-form-approximation}
shows that its continuous--continuous part converges to
\eqref{eq:formal-cc-density} after testing in $(t,x)$.  A mixed term contains
only one continuum integral.  The support of $\chi(\pm(k^2+\kappa_j^2))$ is a
bounded $k$-interval, on which an $L^2$ function is in $L^1$ by
Cauchy--Schwarz.  This justifies the mixed terms as well.

We now disintegrate the temporal Fourier transform.  In the
continuous--continuous term, apply the change of variables
\[
 \delta=k^2-\ell^2,
 \qquad s=k^2+\ell^2.
\]
The shell estimate and Tonelli's theorem show that this change is legitimate
for almost every $\delta$.  Formula \eqref{eq:gap-sum-jacobian} gives
\[
 q_{\alpha\beta,\delta}(s)
 =\frac{
 Q_{\alpha\beta}
 \left(\sqrt{(s+\delta)/2},\sqrt{(s-\delta)/2}\right)
 }{4\sqrt{s^2-\delta^2}}.
 \tag{6.5}\label{eq:q-from-Q}
\]
Define
\[
 G_{\mathrm{cc}}(\delta,x)
 :=\sum_{\alpha,\beta=1}^2
   \int_{|\delta|}^{\infty}
   q_{\alpha\beta,\delta}(s)A_{\alpha\beta}(x,s)\,ds.
\]
For the mixed terms, if $\delta>\kappa_j^2$, put
\[
 k_j^+(\delta)=\sqrt{\delta-\kappa_j^2},
 \qquad
 a^+_{j\alpha}(\delta)
 =\frac{
 g_\alpha(k_j^+(\delta))\overline{c_j}
 -h_\alpha(k_j^+(\delta))\overline{d_j}}
 {2k_j^+(\delta)},
\]
and if $\delta<-\kappa_j^2$, put
\[
 k_j^-(\delta)=\sqrt{-\delta-\kappa_j^2},
 \qquad
 a^-_{j\alpha}(\delta)
 =\frac{
 c_j\overline{g_\alpha(k_j^-(\delta))}
 -d_j\overline{h_\alpha(k_j^-(\delta))}}
 {2k_j^-(\delta)}.
\]
These formulas follow directly from the mixed spectral phases.  The product
of a continuum term with the conjugate of the $j$th bound term has phase
\[
 e^{-it(k^2+\kappa_j^2)},
\]
so it contributes at $\delta=k^2+\kappa_j^2>\kappa_j^2$.  Since
$d\delta=2k\,dk$, its density contains $1/(2k)$.  The conjugate ordering has
phase $e^{it(k^2+\kappa_j^2)}$, contributes at
$\delta=-(k^2+\kappa_j^2)$, and has the same Jacobian.  Set
\begin{align*}
 G_{\mathrm{mix}}(\delta,x)
 :={}&\sum_{j,\alpha}\mathbf 1_{\{\delta>\kappa_j^2\}}
 a^+_{j\alpha}(\delta)
 e_\alpha(x,k_j^+(\delta))\phi_j(x)\\
 &+\sum_{j,\alpha}\mathbf 1_{\{\delta<-\kappa_j^2\}}
 a^-_{j\alpha}(\delta)
 \phi_j(x)\overline{e_\alpha(x,k_j^-(\delta))},
\end{align*}
and $G_{\ac}=G_{\mathrm{cc}}+G_{\mathrm{mix}}$.

With the standard jointly measurable choice of the physical eigenfunctions,
these formulas define strongly measurable $C(K)$-valued functions for every
compact $K$.  To verify continuity in $x$, let $x_n\to x$ in $K$.  Each
product of generalized eigenfunctions converges pointwise, and the local bound
\eqref{eq:scattering-local-bounds} supplies an integrable majorant on every
fixed shell.  Dominated convergence gives convergence in the defining
integrals.

Let $I\subset\R$ be bounded.  By the local scattering bound and the triangle
inequality,
\begin{align*}
 \int_I\|G_{\mathrm{cc}}(\delta,\cdot)\|_{C(K)}\,d\delta
 \lesssim_K\sum_{\alpha,\beta=1}^2
 \iint_{\{k^2-\ell^2\in I\}}
 \bigl(|g_\alpha(k)g_\beta(\ell)|
      +|h_\alpha(k)h_\beta(\ell)|\bigr)\,dk\,d\ell.
\end{align*}
Lemma~\ref{lem:shell} bounds the right side by
\[
 C_{I,K}\sum_{\alpha,\beta=1}^2
 \bigl(\|g_\alpha\|_2\|g_\beta\|_2
       +\|h_\alpha\|_2\|h_\beta\|_2\bigr).
\]
For a positive mixed term, the substitution
$\delta=k^2+\kappa_j^2$ gives
\begin{align*}
 &\int_I\left\|
 \mathbf1_{\{\delta>\kappa_j^2\}}a^+_{j\alpha}(\delta)
 e_\alpha(\cdot,k_j^+(\delta))\phi_j
 \right\|_{C(K)}\,d\delta\\
 &\qquad\lesssim_K (|c_j|+|d_j|)
 \int_{J_{I,j}}(|g_\alpha(k)|+|h_\alpha(k)|)\,dk,
\end{align*}
where $J_{I,j}$ is a bounded $k$-interval.  This is finite by
Cauchy--Schwarz.  The negative mixed term is identical.  Therefore
\[
 \int_I\|G_{\mathrm{mix}}(\delta,\cdot)\|_{C(K)}\,d\delta<\infty.
\]
Consequently
\begin{equation}
 G_{\ac}\in L^1_{\mathrm{loc}}(\R_\delta;C(K))
 \quad\text{for every compact }K\subset\R.
 \tag{6.6}\label{eq:C-K-valued-density}
\end{equation}

Fix $\varphi\in C_c^\infty(\R)$.  The estimates above permit Fubini's theorem,
so the continuous--continuous and mixed parts of the temporal Fourier
transform of \eqref{eq:density-difference}, after pairing in $x$ with
$\varphi$, form the absolutely continuous measure
\[
 \left(\int_\R G_{\ac}(\delta,x)\varphi(x)\,dx\right)d\delta.
\]
A bound--bound product has one fixed temporal frequency
$E_j-E_m$ and hence contributes a multiple of the Dirac mass
$\delta_{E_j-E_m}$.  There are only finitely many such terms.  Thus the
remaining part of the Fourier transform is a finite pure-point measure.
Since \eqref{eq:density-difference} is the zero distribution, its Fourier
transform is the zero distribution.  The Lebesgue decomposition of a locally
finite measure into absolutely continuous and singular parts is unique.
Therefore the displayed scalar density vanishes for almost every $\delta$.  To choose the exceptional set independently of $x$, take the
exhausting compact intervals $K_m=[-m,m]$ and, for each $m$, a countable subset
of $C_c^\infty((-m,m))$ dense in $L^1(K_m)$.  After removing the union of the
corresponding null sets, as well as the null sets on which the chosen
$C(K_m)$-valued representatives are undefined, the integral of
$G_{\ac}(\delta,\cdot)$ against every member of every such dense set is zero.
For a fixed admissible $\delta$, the function
$G_{\ac}(\delta,\cdot)$ is continuous on $K_m$.  If it were nonzero at one
point, continuity would provide an $L^1$ test function for which the integral
is nonzero.  Approximation of that test function by the chosen countable dense
family would give the same conclusion for one member of the family, a
contradiction.  Hence
\begin{equation}
 G_{\ac}(\delta,x)=0
 \quad\text{for every }x\in\R
 \tag{6.7}\label{eq:fixed-gap-common-null-set}
\end{equation}
for almost every $\delta\ne0$.  Intersect this full-measure set with the one
on which every $q_{\alpha\beta,\delta}$ is integrable.

For each remaining $\delta$, equation
\eqref{eq:fixed-gap-common-null-set} is precisely a relation of the form
\eqref{eq:fixed-gap-relation}.  The mixed products are its finitely many
atomic states in the energy-sum variable: their sums are
$s_j^+(\delta)=\delta-2\kappa_j^2$ for $\delta>\kappa_j^2$ and
$s_j^-(\delta)=-\delta-2\kappa_j^2$ for $\delta<-\kappa_j^2$, both strictly
inside the continuum threshold $|s|<|\delta|$.  Proposition~\ref{prop:fixed-gap}
therefore gives $q_{\alpha\beta,\delta}=0$ almost everywhere in $s$ and
$a^\pm_{j\alpha}(\delta)=0$ wherever defined.  This conclusion holds for
almost every $\delta\ne0$.  Under the diffeomorphism
$(k,\ell)\mapsto(\delta,s)$, a null set in the $(\delta,s)$ plane pulls back
to a null set in $(0,\infty)^2$, because the inverse map is locally Lipschitz
away from the boundary.  Formula \eqref{eq:q-from-Q} now gives
\begin{equation}
 Q_{\alpha\beta}(k,\ell)=0
 \quad\text{for every $\alpha,\beta$ and almost every $(k,\ell)\in\R_+^2$}.
 \tag{6.8}\label{eq:ac-kernel-recovered}
\end{equation}
The omitted set $k^2=\ell^2$ is the diagonal and has two-dimensional measure
zero.  Since $k\mapsto\pm(k^2+\kappa_j^2)$ is a diffeomorphism away from
$k=0$, the vanishing of the mixed atomic coefficients gives
\begin{align}
 c_j\overline{g_\alpha(k)}
 &=d_j\overline{h_\alpha(k)},
 \tag{6.9}\label{eq:mixed-one}\\
 g_\alpha(k)\overline{c_j}
 &=h_\alpha(k)\overline{d_j}
 \tag{6.10}\label{eq:mixed-two}
\end{align}
for every $j,\alpha$ and almost every $k>0$.

It remains to recover the finite bound-state block.  Since the part just
analyzed has an $L^1_{\mathrm{loc}}$ temporal-frequency density, it has no
point masses.  The pure-point part of the temporal Fourier transform of
\eqref{eq:density-difference} is therefore the finite sum of the bound--bound
terms.  Dirac masses at distinct gaps are linearly independent as
distributions: one tests with a smooth function supported near one gap and
vanishing near all the others.  Hence, for each fixed gap $\gamma$ and every
spatial test function $\varphi$,
\[
 \int_\R\varphi(x)
 \sum_{E_j-E_m=\gamma}
 \left(c_j\overline{c_m}-d_j\overline{d_m}\right)
 \phi_j(x)\phi_m(x)\,dx=0.
\]
The finite sum in the integrand is continuous.  Its vanishing as a
distribution therefore gives the pointwise identity
\[
 \sum_{E_j-E_m=\gamma}
 \left(c_j\overline{c_m}-d_j\overline{d_m}\right)
 \phi_j(x)\phi_m(x)=0.
 \tag{6.11}\label{eq:bound-gap-relation}
\]
By \eqref{eq:bound-tail-plus}, the summands have nonzero leading terms with
rates
\[
 e^{-(\kappa_j+\kappa_m)x}
 \quad (x\to+\infty).
\]
For a fixed gap
\[
 \gamma=E_j-E_m=\kappa_m^2-\kappa_j^2,
\]
these rates are distinct.  Indeed, knowing $\gamma$ and
$S=\kappa_j+\kappa_m$ determines
\[
 \kappa_m-\kappa_j=\frac{\gamma}{S},
\]
and hence determines the ordered pair $(\kappa_j,\kappa_m)$.  There are finitely many terms in each fixed-gap sum.  Choose the smallest
value of $\kappa_j+\kappa_m$, multiply
\eqref{eq:bound-gap-relation} by $e^{(\kappa_j+\kappa_m)x}$, and let
$x\to+\infty$.  Every faster-decaying term tends to zero, while the selected
term tends to its coefficient times the nonzero product of the leading tail
constants.  Its coefficient must vanish.  Remove that term and repeat.  The
finiteness of the sum yields
\[
 c_j\overline{c_m}=d_j\overline{d_m}
 \quad\text{for all $j,m$}.
 \tag{6.12}\label{eq:bound-kernel-recovered}
\]

\section{Completion of the proof}

Put
\[
 U=\mathcal Fu_0,
 \qquad
 W=\mathcal Fv_0
\]
in
\[
 \HH=\C^N\oplus L^2(\R_+;\C^2).
\]
For $Y\in\HH$, define the rank-one operator generated by $U$ by
\[
 (U\otimes\overline U)Y=\langle Y,U\rangle_{\HH}U.
\]
Its kernel has four blocks: bound--bound, bound--continuous,
continuous--bound, and continuous--continuous.  Equation
\eqref{eq:bound-kernel-recovered} identifies the first block;
\eqref{eq:mixed-one}--\eqref{eq:mixed-two} identify the two off-diagonal
blocks; and \eqref{eq:ac-kernel-recovered} identifies the final block.
Therefore
\[
 U\otimes\overline U=W\otimes\overline W
 \tag{7.1}\label{eq:rank-one-equality}
\]
as operators on $\HH$.

If $U=0$, the left side is zero.  Taking the inner product of the right side
with $W$ gives $\|W\|_{\HH}^4=0$, so $W=0$; take $\zeta=1$.  Suppose
$U\ne0$.  The range of $U\otimes\overline U$ is $\operatorname{span}\{U\}$.
Equality \eqref{eq:rank-one-equality} shows that its range also equals
$\operatorname{span}\{W\}$, so $W=\zeta U$ for a nonzero scalar $\zeta$.
Substitution in \eqref{eq:rank-one-equality} gives
\[
 U\otimes\overline U=|\zeta|^2U\otimes\overline U.
\]
Since the operator on the left is nonzero, $|\zeta|=1$.  In either case,
since $\mathcal F$ is unitary,
\[
 v_0=\zeta u_0.
\]
This proves Theorem~\ref{thm:main}.

\part{Further results and remarks}\label{part:further-remarks}

\section{One-sided observation for compactly supported potentials}

We prove Theorem~\ref{thm:half-line}.  We first isolate the Fourier statement
used in the proof.  For $F\in L^2(\R)$ set
\[
 \mathcal E_tF=\int_\R F(k)e^{iky-ik^2t}\,dk
\]
as an $L^2_y$ inverse Fourier transform.

\begin{lemma}[Free half-line model]
\label{lem:free-half-line}
Let $0<\kappa_1<\cdots<\kappa_N$, let $a_j,b_j\in\C$, and let
$F,G\in L^2(\R)$.  Suppose
\begin{align*}
 U(t,y)&=\sum_{j=1}^Na_je^{-\kappa_jy}e^{i\kappa_j^2t}
          +\mathcal E_tF(y),\\
 W(t,y)&=\sum_{j=1}^Nb_je^{-\kappa_jy}e^{i\kappa_j^2t}
          +\mathcal E_tG(y)
\end{align*}
in $L^2((0,\infty)_y)$ for every $t$, and assume
$|U(t,y)|=|W(t,y)|$ for almost every $(t,y)\in\R\times(0,\infty)$.
Then there is $\zeta\in\T$ such that
\[
 G=\zeta F\quad\text{a.e. on }\R,
 \qquad b_j=\zeta a_j\quad(1\le j\le N).
\]
The assertion includes $N=0$.
\end{lemma}

\begin{proof}
Plancherel's theorem gives
$\|\mathcal E_tF\|_{L^2(\R)}=\sqrt{2\pi}\|F\|_2$, independently of $t$.
Moreover, dominated convergence in Fourier space shows that
$t\mapsto\mathcal E_tF$ is continuous into $L^2(\R)$.  The finite bound-state
sum is also continuous into $L^2(0,\infty)$.  Hence
$t\mapsto U(t),W(t)$ are continuous into $L^2(0,\infty)$.  The elementary
estimate
\[
 \||a|^2-|b|^2\|_1\le\|a-b\|_2(\|a\|_2+\|b\|_2)
\]
shows that $t\mapsto |U(t)|^2-|W(t)|^2$ is continuous into $L^1(0,\infty)$.
It vanishes for almost every $t$ by hypothesis, and therefore vanishes for
every $t$.

We first record a Cauchy-transform fact.  If $J\subset\R$ is a bounded open
interval, $m\in L^1(J)$, and
\[
 C_m(z)=\int_J\frac{m(\eta)}{z-i\eta}\,d\eta,
 \qquad z\in\C\setminus iJ,
\]
extends holomorphically through $iJ$, then $m=0$ almost everywhere on $J$.
Indeed,
\[
 H(w):=iC_m(iw)=\int_J\frac{m(\eta)}{w-\eta}\,d\eta
\]
extends through $J$.  Let $K\Subset J$ and
$\phi\in C_c^\infty(K)$.  For $\varepsilon>0$ small enough, Fubini's theorem
and
\[
 \frac1{x+i\varepsilon-\eta}-\frac1{x-i\varepsilon-\eta}
 =\frac{-2i\varepsilon}{(x-\eta)^2+\varepsilon^2}
\]
give
\begin{align*}
 &\int_K\bigl(H(x+i\varepsilon)-H(x-i\varepsilon)\bigr)\phi(x)\,dx\\
 &\qquad=-2\pi i\int_Jm(\eta)(P_\varepsilon*\phi)(\eta)\,d\eta,
 \qquad
 P_\varepsilon(x)=\frac1\pi\frac{\varepsilon}{x^2+\varepsilon^2}.
\end{align*}
Because the extended function $H$ is uniformly continuous on a compact
neighborhood of $K$, the left side tends to zero.  The Poisson kernels form
an approximate identity, so $P_\varepsilon*\phi\to\phi$ uniformly.  Since
$m\in L^1(J)$, the right side tends to
$-2\pi i\int_Jm(\eta)\phi(\eta)\,d\eta$.  Thus $m$ vanishes as a
distribution on every $K\Subset J$, and hence almost everywhere on $J$.

Fix $\chi\in C_c^\infty(\R)$ and define
\[
 \check\chi(t)=\frac1{2\pi}\int_\R e^{-i\omega t}\chi(\omega)\,d\omega.
\]
The Fourier inversion formula gives
\[
 \int_\R \check\chi(t)e^{i\omega t}\,dt=\chi(\omega).
 \tag{8.0}\label{eq:half-line-frequency-selector}
\]
For $\Re z>0$, multiply $|U|^2-|W|^2$ by
$\check\chi(t)e^{-zy}$ and integrate in $(t,y)$.  This integral is absolutely
convergent.  Indeed, if $a=\Re z>0$, then
\[
 \int_0^\infty e^{-ay}|\mathcal E_tF(y)|^2\,dy
 \le\|\mathcal E_tF\|_2^2=2\pi\|F\|_2^2
\]
uniformly in $t$, and the same estimate holds for $G$.  Each bound-state sum
has a uniformly bounded exponentially weighted $L^2_y$ norm.  Since
$\check\chi\in L^1(\R)$, integration in $t$ is legitimate.

For Schwartz $F,G$, expand the density difference into free--free,
bound--free, free--bound, and bound--bound terms.  For example, the
free--free term is
\[
 \iint_{\R^2}q(k,\ell)
 e^{i(k-\ell)y}e^{i(\ell^2-k^2)t}\,dk\,d\ell.
\]
Using \eqref{eq:half-line-frequency-selector} and
\[
 \int_0^\infty e^{-zy}e^{i(k-\ell)y}\,dy
 =\frac1{z-i(k-\ell)}
\]
produces $\Phi_\chi$.  A bound--free term has spatial factor
$e^{-(\kappa_j+ik)y}$ and temporal frequency $k^2+\kappa_j^2$, producing
$\Psi^+_{j,\chi}$.  The opposite ordering produces
$\Psi^-_{j,\chi}$, and two bound states produce $R_\chi$.  Thus
\begin{equation}
 0=\Phi_\chi(z)+\sum_{j=1}^N\bigl(\Psi^+_{j,\chi}(z)
       +\Psi^-_{j,\chi}(z)\bigr)+R_\chi(z),
 \tag{8.1}\label{eq:half-line-transform}
\end{equation}
where
\begin{align*}
 q(k,\ell)&=F(k)\overline{F(\ell)}-G(k)\overline{G(\ell)},\\
 r_j^+(k)&=a_j\overline{F(k)}-b_j\overline{G(k)},\\
 r_j^-(k)&=\overline{a_j}F(k)-\overline{b_j}G(k),\\
 d_{jm}&=a_j\overline{a_m}-b_j\overline{b_m},
\end{align*}
and
\begin{align*}
 \Phi_\chi(z)
 &=\iint_{\R^2}q(k,\ell)
   \frac{\chi(\ell^2-k^2)}{z-i(k-\ell)}\,dk\,d\ell,\\
 \Psi^+_{j,\chi}(z)
 &=\int_\R r_j^+(k)
   \frac{\chi(k^2+\kappa_j^2)}{z+\kappa_j+ik}\,dk,\\
 \Psi^-_{j,\chi}(z)
 &=\int_\R r_j^-(k)
   \frac{\chi(-(k^2+\kappa_j^2))}{z+\kappa_j-ik}\,dk,\\
 R_\chi(z)
 &=\sum_{j,m=1}^Nd_{jm}
   \frac{\chi(\kappa_j^2-\kappa_m^2)}{z+\kappa_j+\kappa_m}.
\end{align*}
To justify \eqref{eq:half-line-transform} for arbitrary $L^2$ data, choose
Schwartz functions $F_n,G_n$ converging to $F,G$ in $L^2$.  Let $U_n,W_n$ be
the corresponding functions with the bound-state coefficients unchanged.
Plancherel gives convergence uniformly in $t$ in $L^2_y$, and the inequality
\[
 \||p|^2-|q|^2\|_1\le\|p-q\|_2(\|p\|_2+\|q\|_2)
\]
shows that the transformed density of $U_n,W_n$ converges to that of $U,W$.
The factor $e^{-\Re(z)y}$ only improves this estimate, and
$\check\chi\in L^1_t$ permits integration in time.

It remains to pass to the limit on the spectral side.  If
$\supp\chi\subset[-M,M]$, the kernel
$\1_{\{|\ell^2-k^2|\le M\}}$ has uniformly bounded integrals in each
variable.  Schur's test therefore gives
\[
 \iint |F(k)F(\ell)|\,|\chi(\ell^2-k^2)|\,dk\,d\ell
 \lesssim_{\chi}\|F\|_2^2,
\]
and the corresponding bilinear estimate controls the difference between the
kernels formed from $F_n$ and $F$.  The same statement holds for $G$.
The mixed factors have compact $k$-support and are controlled by
Cauchy--Schwarz.  Thus every term in \eqref{eq:half-line-transform} converges
to the displayed formula.

We next separate the free--free term from the others by analytic
continuation.  The numerator of $\Psi^+_{j,\chi}$ has compact $k$-support.
If $|\Re z|<\kappa_j/2$, then
\[
 |z+\kappa_j+ik|\ge \Re(z+\kappa_j)>\kappa_j/2.
\]
The integrand and all its $z$-derivatives are therefore dominated by
integrable functions on compact substrips.  Thus $\Psi^+_{j,\chi}$ is
holomorphic there.  The same argument applies to $\Psi^-_{j,\chi}$.
The poles of $R_\chi$ are the negative real numbers
$-(\kappa_j+\kappa_m)$, so $R_\chi$ is holomorphic in a strip around the
imaginary axis.  Equation \eqref{eq:half-line-transform} now defines a
holomorphic continuation of $\Phi_\chi$ through that axis.  Put
\[
 M_\chi(\xi)=\int_\R q(k,k-\xi)
                 \chi(\xi^2-2k\xi)\,dk.
\]
The preceding Schur estimate and Fubini's theorem give $M_\chi\in L^1(\R)$ and
\[
 \Phi_\chi(z)=\int_\R\frac{M_\chi(\xi)}{z-i\xi}\,d\xi.
\]
For a bounded interval $J$, split this integral into its parts over $J$ and
$\R\setminus J$.  The latter is holomorphic near $iJ$, so the
Cauchy-transform fact applies to the former.  An exhaustion by bounded
intervals gives $M_\chi=0$ almost everywhere for every $\chi$.

This implies $q=0$ almost everywhere on $\R^2$.  We include the measure-theoretic
step.  Away from $\xi=0$, the change of variables
\[
 \xi=k-\ell,
 \qquad \omega=\ell^2-k^2
\]
has Jacobian $2|\xi|$ and inverse
\[
 k=\frac{\xi^2-\omega}{2\xi},
 \qquad
 \ell=-\frac{\xi^2+\omega}{2\xi}.
\]
Let $I\Subset\R\setminus\{0\}$ and $J\Subset\R$.  For
$\theta\in C_c^\infty(I)$ and $\eta\in C_c^\infty(J)$, the identity
$M_\eta=0$ gives
\[
 0=\int_I\theta(\xi)\int_\R q(k,k-\xi)
 \eta(\xi^2-2k\xi)\,dk\,d\xi.
\]
On the support of the integrand, use $(\xi,\omega)$ as variables.  Since
$dk\,d\xi=(2|\xi|)^{-1}d\omega\,d\xi$, this becomes
\[
 \int_{I\times J}
 \frac{q(k(\xi,\omega),\ell(\xi,\omega))}{2|\xi|}
 \theta(\xi)\eta(\omega)\,d\xi\,d\omega=0.
\]
The pullback of $q$ belongs to $L^2(I\times J)$, and $1/(2|\xi|)$ is bounded
above and below there.  Finite sums of the tensors
$\theta(\xi)\eta(\omega)$ are dense in $L^2(I\times J)$.  The weighted
pullback, and hence the pullback itself, vanishes almost everywhere.  Choose
countable families of such intervals exhausting
$(\R\setminus\{0\})\times\R$.  The inverse change of variables then gives
$q(k,\ell)=0$ away from the diagonal $k=\ell$.  The diagonal has
Lebesgue measure zero, so $q=0$ almost everywhere on $\R^2$.

Consequently the Hilbert--Schmidt kernels of the two rank-one operators
coincide:
\[
 F(k)\overline{F(\ell)}=G(k)\overline{G(\ell)}
 \quad\text{for almost every }(k,\ell).
\]
If $F=0$, Fubini's theorem shows that $G(k)\overline{G(\ell)}=0$ almost
everywhere, hence $G=0$.  If $F\ne0$, the two nonzero rank-one operators have
the same range, so $G=\zeta F$ for some nonzero scalar $\zeta$.  Substitution
in the kernel identity gives $|\zeta|^2=1$.  Thus $|\zeta|=1$.

It remains to identify the finite sums.  First suppose $F\ne0$.  Replace $W$
by $\overline\zeta W$; then $G=F$.  We must prove $b_j=a_j$ in this
normalization.  Set
\[
 \Lambda=\{\kappa_j^2-\kappa_m^2:1\le j,m\le N\}.
\]
Choose $\chi\in C_c^\infty((0,\infty)\setminus\Lambda)$.  In
\eqref{eq:half-line-transform}, the free--free, negative-frequency mixed, and
bound--bound terms vanish.  Hence
\[
 \sum_{j=1}^NC_{j,\chi}(z)=0,
 \qquad
 C_{j,\chi}(z)=\int_\R
 \frac{(a_j-b_j)\overline{F(k)}\chi(k^2+\kappa_j^2)}
      {z+\kappa_j+ik}\,dk.
 \tag{8.2}\label{eq:bound-induction}
\]
The function $C_{j,\chi}$ is holomorphic on the half-plane
$\Re z>-\kappa_j$.  The sum in \eqref{eq:bound-induction} vanishes on
$\Re z>0$, so the identity theorem makes it vanish wherever all terms are
holomorphic.

We begin with $j=1$.  Every $C_{m,\chi}$ with $m>1$ is holomorphic on a
neighborhood of the line $\Re z=-\kappa_1$, because
$-\kappa_1>-\kappa_m$.  On the right of that line,
\[
 C_{1,\chi}(z)=-\sum_{m=2}^NC_{m,\chi}(z).
\]
The right side extends through the line, and therefore so does
$C_{1,\chi}$.  Put $w=z+\kappa_1$ and change $k$ to $-\eta$.  Then
$C_{1,\chi}(w-\kappa_1)$ is the Cauchy transform of the compactly supported
$L^1$ function
\[
 \eta\longmapsto
 (a_1-b_1)\overline{F(-\eta)}
 \chi(\eta^2+\kappa_1^2).
\]
The Cauchy-transform fact gives
\[
 (a_1-b_1)\overline{F(k)}\chi(k^2+\kappa_1^2)=0
 \quad\text{for almost every }k.
\]
Thus $C_{1,\chi}=0$.  Remove it from \eqref{eq:bound-induction}.  The same
argument at $\Re z=-\kappa_2$ gives the result for $j=2$.  Induction gives
\[
 (a_j-b_j)\overline{F(k)}\chi(k^2+\kappa_j^2)=0
 \quad\text{for almost every }k
\]
for every $j$.  Varying $\chi$ shows that
$(a_j-b_j)F(k)=0$ outside the finite set
$\{k:k^2+\kappa_j^2\in\Lambda\}$.  Since $F\ne0$, one has $a_j=b_j$.

If $F=G=0$, expansion of the modulus identity gives
\[
 \sum_{j,m=1}^Nd_{jm}e^{-(\kappa_j+\kappa_m)y}
 e^{i(\kappa_j^2-\kappa_m^2)t}=0.
\]
The left side is continuous, so the almost-everywhere identity holds for all
$(t,y)$.  Fix $y>0$ and group terms having the same temporal frequency
$\omega=\kappa_j^2-\kappa_m^2$.  Distinct exponentials $e^{i\omega t}$ are
linearly independent: evaluating the first $M$ derivatives at one time gives
a Vandermonde system.  Hence, for every $\omega$,
\[
 \sum_{\kappa_j^2-\kappa_m^2=\omega}
 d_{jm}e^{-(\kappa_j+\kappa_m)y}=0
 \quad(y>0).
\]
Within one frequency class, the numbers $\kappa_j+\kappa_m$ are distinct.
Indeed, the difference of squares and the sum determine the ordered pair:
\[
 \kappa_j-\kappa_m=
 \frac{\omega}{\kappa_j+\kappa_m}.
\]
A second Vandermonde argument, now in $y$, gives $d_{jm}=0$ for all $j,m$.  Hence
$(a_1,\ldots,a_N)\otimes(\overline a_1,\ldots,\overline a_N)$ equals the
corresponding rank-one matrix for $b$, and $b=\zeta a$ for some $\zeta\in\T$.
This completes the proof.
\end{proof}

\begin{proof}[Proof of Theorem~\ref{thm:half-line}]
Translation reduces the proof to $R=0$.  Standard one-dimensional scattering
theory gives finitely many simple negative eigenvalues
$-\kappa_1^2,\ldots,-\kappa_N^2$, real normalized eigenfunctions $\phi_j$,
and two continuum channels $e_1(\cdot,k),e_2(\cdot,k)$, $k>0$.  Since $V=0$
on $(0,\infty)$, square integrability excludes the growing exponential and
therefore
\[
 \phi_j(x)=\alpha_je^{-\kappa_jx},\qquad \alpha_j\ne0.
\]
The coefficient cannot vanish, since the eigenfunction and its derivative
would then vanish on $(0,\infty)$ and ODE uniqueness would give
$\phi_j\equiv0$.
and the continuum functions may be normalized so that
\[
 e_1(x,k)=T(k)e^{ikx},
 \qquad
 e_2(x,k)=e^{-ikx}+R_2(k)e^{ikx},
 \qquad x>0.
\]
The coefficients $T,R_2$ are bounded, and $T(k)\ne0$ for $k>0$.  Indeed,
$T(k_0)=0$ would make $e_1(\cdot,k_0)$ and its derivative vanish on
$(0,\infty)$; ODE uniqueness would then give $e_1\equiv0$, contradicting its
nonzero incoming asymptotic at $-\infty$.

Write the spectral coefficients of $u_0$ as
$c_j;g_1,g_2$ and those of $v_0$ as $d_j;h_1,h_2$.  Suppressing one common
normalization constant, the spectral expansion on $x>0$ is
\begin{align*}
 e^{-itH}u_0(x)
 ={}&\sum_{j=1}^Nc_j\alpha_je^{-\kappa_jx}e^{i\kappa_j^2t}\\
 &+\int_0^\infty e^{-ik^2t}
 \bigl[T(k)g_1(k)e^{ikx}
 +g_2(k)e^{-ikx}+R_2(k)g_2(k)e^{ikx}\bigr]dk.
\end{align*}
Combine the two coefficients of $e^{ikx}$ and change $k$ to $-k$ in the
coefficient of $e^{-ikx}$.  The continuum integral becomes
$\int_\R F_u(k)e^{ikx-ik^2t}\,dk$, where
\[
 F_u(k)=
 \begin{cases}
 T(k)g_1(k)+R_2(k)g_2(k),&k>0,\\
 g_2(-k),&k<0.
 \end{cases}
\]
Define $F_v$ analogously.  The boundedness of $T$ and $R_2$ implies
$F_u,F_v\in L^2(\R)$.  Thus the restrictions of the two evolutions to the free
half-line have exactly the form in Lemma~\ref{lem:free-half-line}.
The half-line modulus identity and Lemma~\ref{lem:free-half-line} give
\[
 F_v=\zeta F_u,
 \qquad d_j\alpha_j=\zeta c_j\alpha_j.
\]
Thus $d_j=\zeta c_j$.  On the negative frequency half-line,
$F_v=\zeta F_u$ gives $h_2=\zeta g_2$.  On the positive frequency half-line it
then gives
\[
 T(k)(h_1(k)-\zeta g_1(k))=0
 \quad\text{a.e. }k>0.
\]
Since $T(k)\ne0$, $h_1=\zeta g_1$.  All spectral coefficients agree up to
$\zeta$, and unitarity of the spectral transform gives $v_0=\zeta u_0$.
\end{proof}

\subsection*{Scope of the decay hypothesis}

The product system only requires $V\in L^1_{\mathrm{loc}}$, and the
negative-energy argument uses integrable tails.  The first moment in
Theorem~\ref{thm:main} enters through the scattering representation, its
uniform threshold bounds, and the finiteness of the negative spectrum.  One
could replace $V\in L^1_1$ by $V\in L^1$ only after imposing and verifying the
needed threshold and negative-spectrum hypotheses.  No such extension is
claimed here.

\section{Phase retrieval from Masuda's theorem}

Let $V:\R\to\R$ be measurable.  We impose the assumptions in
Masuda's theorem~\cite{Masuda}:
\begin{enumerate}
\item $V\in L^2_{\mathrm{loc}}(\R)$;
\item there is a null set $Q\subset\R$ such that $\R\setminus Q$ is connected
and open and $V$ is locally essentially bounded on $\R\setminus Q$;
\item the operator
\[
 H_0=-\partial_x^2+V,
 \qquad D(H_0)=C_c^\infty(\R),
\]
is essentially self-adjoint on $L^2(\R)$;
\item the spectrum of its self-adjoint closure $H$ is bounded from below.
\end{enumerate}
Masuda also permits a real $C^2$ magnetic potential.  We only use the
nonmagnetic case.

A Masuda weak solution of
\begin{equation}
 if_t+f_{xx}=Vf
 \tag{9.1}\label{eq:schrodinger}
\end{equation}
is a function $f\in C(\R;L^2(\R))$ satisfying
\eqref{eq:schrodinger} in distributions.  Thus, for every
$\varphi\in C_c^\infty(\R^2)$,
\[
 \iint_{\R^2}f(t,x)
 \bigl(-i\varphi_t(t,x)+\varphi_{xx}(t,x)-V(x)\varphi(t,x)\bigr)
 \,dx\,dt=0,
\]
with the usual complex-bilinear distributional pairing.  A preliminary lemma
in Masuda's proof identifies every such weak solution with the unitary orbit
\[
 f(t)=e^{-i(t-t_0)H}f(t_0).
\]
The same time-mollification argument on a compact subinterval gives local
uniqueness; hence a weak solution defined only on an open time interval has a
unique global unitary extension from any one of its time slices.  Under the
assumptions above, Masuda proves the following statement.

\begin{proposition}[Masuda]
\label{prop:masuda}
If a Masuda weak solution of \eqref{eq:schrodinger} vanishes on a nonempty
open subset of $\R_t\times\R_x$, then it vanishes identically.
\end{proposition}

We now prove Theorem~\ref{thm:phase-retrieval}.  Thus $I\subset\R$ is a
nonempty open interval, $u,v$ are Masuda weak solutions, and
\[
 u,v\in C(I;H^1(\R)).
 \tag{9.2}\label{eq:H1-regularity}
\]
The density hypothesis is
\begin{equation}
 |u(t,x)|=|v(t,x)|
 \quad\text{for almost every }(t,x)\in I\times\R.
 \tag{9.3}\label{eq:modulus}
\end{equation}
Equality of the densities first determines the difference of the currents.
The one-dimensional integrability of that difference is used below.

\section{Density and current}

For $f\in H^1(\R)$ define
\[
 \rho_f=|f|^2,
 \qquad
 j_f=\operatorname{Im}(\overline f f_x).
\]
Then $j_f\in L^1(\R)$ and
\[
 \|j_f\|_1\le \|f\|_2\|f_x\|_2.
\]

\begin{lemma}[Recovery of the current]
\label{lem:current}
Under the hypotheses of Theorem~\ref{thm:phase-retrieval},
\begin{equation}
 j_u=j_v
 \quad\text{almost everywhere on }I\times\R.
 \tag{10.1}\label{eq:equal-current}
\end{equation}
Consequently,
\begin{equation}
 \overline u\,u_x=\overline v\,v_x
 \quad\text{almost everywhere on }I\times\R.
 \tag{10.2}\label{eq:equal-logarithmic-derivative}
\end{equation}
\end{lemma}

\begin{proof}
Every solution satisfying \eqref{eq:H1-regularity} obeys the continuity
equation
\begin{equation}
 \partial_t|f|^2+2\partial_x
 \operatorname{Im}(\overline f f_x)=0
 \quad\text{in }\D'(I\times\R).
 \tag{10.3}\label{eq:continuity}
\end{equation}
For a smooth solution, solve \eqref{eq:schrodinger} for the time derivative:
$f_t=if_{xx}-iVf$.  Since $V$ is real,
\begin{align*}
 \partial_t|f|^2
 &=f_t\overline f+f\overline{f_t}\\
 &=i(f_{xx}\overline f-f\overline{f_{xx}})\\
 &=-2\partial_x\operatorname{Im}(\overline f f_x).
\end{align*}
For a solution satisfying \eqref{eq:H1-regularity}, regularize the equation in
$(t,x)$ and repeat this calculation.  On compact spatial intervals,
$V\in L^2$, while the one-dimensional embedding
$H^1\subset L^\infty$ gives $Vf\in L^2_{\mathrm{loc}}$.  Thus the potential
terms converge and cancel.  The products containing $f_x$ converge in
$L^1$ by Cauchy--Schwarz and \eqref{eq:H1-regularity}.  Passing to the limit
proves \eqref{eq:continuity}.

Subtract \eqref{eq:continuity} for $u$ and $v$.  By
\eqref{eq:modulus},
\[
 \partial_x(j_u-j_v)=0
 \quad\text{in }\D'(I\times\R).
 \tag{10.4}\label{eq:current-constant}
\]
We show that the constant in $x$ is zero.  For $\eta\in C_c^\infty(I)$ put
\[
 J_\eta(x)=\int_I\eta(t)(j_u-j_v)(t,x)\,dt.
\]
The estimate
\[
 \|j_f(t)\|_1\le\|f(t)\|_2\|f_x(t)\|_2
\]
and the $C(I;H^1)$ assumption give $J_\eta\in L^1(\R)$.  Testing
\eqref{eq:current-constant} against $\eta(t)\phi(x)$ shows that
$J_\eta'=0$ in $\D'(\R)$.  A distribution on $\R$ with zero derivative is a
constant distribution: for example, it annihilates every test function of
integral zero, each of which is the derivative of a compactly supported test
function.  Since $J_\eta$ is represented by an $L^1$ function, this constant
must be zero.  Thus
\[
 \int_I\eta(t)(j_u-j_v)(t,x)\,dt=0
\]
for almost every $x$ and every temporal test function $\eta$.  Hence
\eqref{eq:equal-current} holds.

Finally,
\[
 \operatorname{Re}(\overline f f_x)
 =\frac12\partial_x|f|^2
\]
in distributions.  The real parts of $\overline u u_x$ and
$\overline v v_x$ agree by \eqref{eq:modulus}, and their imaginary parts agree
by \eqref{eq:equal-current}.  Both products belong locally to $L^1$, so the
distributional identity is the almost-everywhere identity
\eqref{eq:equal-logarithmic-derivative}.
\end{proof}

\section{Local constancy of the relative phase}

Every $H^1(\R)$ function has a continuous representative, since
\[
 |f(x)-f(y)|\le |x-y|^{1/2}\|f_x\|_2.
\]
Together with the $L^2$ condition, this representative tends to zero at
infinity.  The embedding $H^1(\R)\hookrightarrow C_0(\R)$ is continuous.
Thus \eqref{eq:H1-regularity} gives jointly continuous representatives of
$u$ and $v$: continuity in $t$ holds uniformly in $x$, and continuity in $x$
holds on every time slice.  Set
\[
 \Omega=\{(t,x)\in I\times\R:u(t,x)\ne0\}.
\]
By \eqref{eq:modulus}, $v$ has the same nonzero set.  The set $\Omega$ is
open.

\begin{lemma}[Relative phase]
\label{lem:relative-phase}
On every connected component of $\Omega$ there is a constant $\zeta\in\T$
such that $v=\zeta u$.
\end{lemma}

\begin{proof}
Fix an open rectangle $J\times O$ whose closure is contained in $\Omega$.
There is a constant $c_0>0$ such that
\[
 |u(t,x)|\ge c_0
 \quad ((t,x)\in J\times O).
\]
Define $q=v/u$ there.  The map $z\mapsto1/z$ is smooth on
$\{|z|\ge c_0\}$, so the Sobolev chain and product rules give
$q\in C(J;H^1(O))$.  Equation \eqref{eq:modulus} and continuity imply
$|q|=1$ everywhere on the rectangle.  For almost
every $t\in J$, the quotient rule in $H^1(O)$ and
\eqref{eq:equal-logarithmic-derivative} give
\begin{align*}
 0
 &=\overline v v_x-\overline u u_x\\
 &=|u|^2\overline q\,q_x
 \quad\text{almost everywhere on }O.
\end{align*}
Hence $q(t,\cdot)$ is constant on $O$ for almost every $t$.  Since
$q\in C(J;H^1(O))\subset C(J\times O)$, approximation by such times gives the
same conclusion for every $t\in J$.  We may therefore write
\[
 v(t,x)=c(t)u(t,x)
 \quad ((t,x)\in J\times O),
 \qquad |c(t)|=1.
 \tag{11.1}\label{eq:local-time-phase}
\]

It remains to prove that $c$ is independent of $t$.  Choose a nonnegative,
nonzero function $\chi\in C_c^\infty(O)$ and set
\[
 A(t)=\int_\R\chi(x)\overline{u(t,x)}v(t,x)\,dx,
 \qquad
 B(t)=\int_\R\chi(x)|u(t,x)|^2\,dx.
\]
The lower bound for $|u|$ on $\supp\chi$ gives $B(t)>0$ on $J$, and
\eqref{eq:local-time-phase} gives
\[
 A(t)=c(t)B(t).
 \tag{11.2}\label{eq:A-cB}
\]
We compute their derivatives.  Formally,
\[
 u_t=iu_{xx}-iVu,
 \qquad
 v_t=iv_{xx}-iVv.
\]
Therefore
\begin{align*}
 A'
 &=\int\chi(\overline{u_t}v+\overline u v_t)\\
 &=i\int\chi(\overline u v_{xx}-\overline{u_{xx}}v),
\end{align*}
because the two potential terms cancel.  Integrating each second derivative
once transfers the remaining derivative to $\chi$ and gives, for almost every
$t\in J$,
\begin{align}
 A'(t)
 &=i\int_\R\chi'(x)
   \bigl(\overline{u_x(t,x)}v(t,x)
         -\overline{u(t,x)}v_x(t,x)\bigr)\,dx,
 \tag{11.3}\label{eq:A-prime}\\
 B'(t)
 &=i\int_\R\chi'(x)
   \bigl(\overline{u_x(t,x)}u(t,x)
         -\overline{u(t,x)}u_x(t,x)\bigr)\,dx.
 \tag{11.4}\label{eq:B-prime}
\end{align}
The same calculation with $v=u$ gives \eqref{eq:B-prime}.  All displayed
integrands belong to $L^1$ by Cauchy--Schwarz.  For the weak solutions, perform
the calculation after local regularization.  The terms
$V\overline uv$ are locally integrable because
$V\in L^2_{\mathrm{loc}}$ and $u,v$ are locally bounded.  The formulas show
that the distributional derivatives $A'$ and $B'$ belong to
$L^1_{\mathrm{loc}}(J)$; hence $A$ and $B$ have locally absolutely continuous
representatives.

On $\supp\chi$, equation \eqref{eq:local-time-phase} and its weak spatial
derivative give $v=cu$ and $v_x=cu_x$.  Therefore
\[
 A'(t)=c(t)B'(t)
 \quad\text{for almost every }t\in J.
\]
Since $B>0$, equation \eqref{eq:A-cB} shows that $c=A/B$ is locally
absolutely continuous.  Differentiating \eqref{eq:A-cB} gives
\[
 c'(t)B(t)=A'(t)-c(t)B'(t)=0
\]
for almost every $t$.  Hence $c$ is constant on $J$.

Every point of $\Omega$ lies in a rectangle of the type just considered, so
$q$ is locally constant on $\Omega$.  Let $\Omega_0$ be a connected
component and fix one value $\zeta$ attained there.  The set
$\{(t,x)\in\Omega_0:q(t,x)=\zeta\}$ is nonempty, open, and closed relative to
$\Omega_0$.  Connectedness implies that it is all of $\Omega_0$.  This proves
the lemma.
\end{proof}

\section{Proof of phase retrieval}

\begin{proof}[Proof of Theorem~\ref{thm:phase-retrieval}]
If $u$ vanishes on $I\times\R$, Proposition~\ref{prop:masuda} gives
$u\equiv0$.  Equation \eqref{eq:modulus} then shows that $v$ vanishes on the
same nonempty open set, so Proposition~\ref{prop:masuda} also gives
$v\equiv0$.  The conclusion holds with $\zeta=1$.

Suppose that $\Omega$ is nonempty.  Choose one connected component
$\Omega_0$.  By Lemma~\ref{lem:relative-phase}, there is a $\zeta\in\T$ such
that
\[
 v=\zeta u\quad\text{on }\Omega_0.
\]
The equation is linear and $u,v$ have the same potential.  Hence the
difference
\[
 w=v-\zeta u
\]
belongs to $C(\R;L^2)$ and is a Masuda weak solution of
\eqref{eq:schrodinger}.  It vanishes on the nonempty open subset
$\Omega_0$ of spacetime.  Proposition
\ref{prop:masuda} gives $w\equiv0$ on $\R^2$.  Thus
$v(t)=\zeta u(t)$ in $L^2(\R)$ for every $t\in\R$.
\end{proof}

\section{The \texorpdfstring{$L^2$}{L2} endpoint in Masuda's class}
\label{sec:masuda-L2}

Masuda's unique-continuation theorem is already formulated for
$C(\R;L^2)$ weak solutions.  It does not, by itself, remove
\eqref{eq:H1-regularity} from Theorem~\ref{thm:phase-retrieval}.  The missing
step occurs before unique continuation: one must first construct a constant
$\zeta$ for which $v-\zeta u$ vanishes on a spacetime-open set.

There is also a minor simplification peculiar to one dimension.  In Masuda's
original hypothesis, $\R\setminus Q$ is connected and open and has full
measure.  Such a subset of $\R$ must be all of $\R$.  Thus the literal
one-dimensional assumption forces $Q=\varnothing$ and
$V\in L^\infty_{\mathrm{loc}}(\R)$.  This removes potential-product
singularities, but it does not define the kinetic products used below.

\subsection*{The current is not an \texorpdfstring{$L^2$}{L2} operation}

For $f\in L^2$, one only has $f_x\in H^{-1}$.  The product
$\overline f f_x$ is not defined by Sobolev duality, and it has no extension
to all of $L^2$ which is continuous under smooth $L^2$ approximation and
agrees with the classical product.

We give an explicit obstruction.  Choose a nonzero real Schwartz function
$\varphi$ whose Fourier transform has sufficiently small compact support, and
put
\[
 \lambda_n=3^n,
 \qquad
 a_n=(n\lambda_n)^{-1/2},
 \qquad
 f_N(x)=\varphi(x)\sum_{n=1}^Na_ne^{i\lambda_nx}.
 \tag{13.1}\label{eq:lacunary-current-example}
\]
The shifted Fourier supports are disjoint.  Hence $(f_N)$ converges in
$L^2(\R)$, since
\[
 \sum_{n=1}^\infty a_n^2
 =\sum_{n=1}^\infty\frac1{n3^n}<\infty.
\]
Let $\chi\in C_c^\infty(\R)$ be nonnegative and satisfy
$A_\chi:=\int\chi|\varphi|^2>0$.  The diagonal frequency terms give
\begin{align}
 \int_\R\chi\,\Im(\overline{f_N}(f_N)_x)
 &=A_\chi\sum_{n=1}^N\lambda_na_n^2+R_N \notag\\
 &=A_\chi\sum_{n=1}^N\frac1n+R_N.
 \tag{13.2}\label{eq:divergent-current}
\end{align}
The remainder $(R_N)$ is bounded.  Indeed, every off-diagonal term contains
the Fourier transform of either $\chi|\varphi|^2$ or
$\chi\varphi\varphi'$, evaluated at $\lambda_m-\lambda_n$.  For every
integer $M\ge1$, both transforms are bounded by
$C_M(1+|\lambda_m-\lambda_n|)^{-M}$.  Lacunarity gives
\[
 |\lambda_m-\lambda_n|
 \ge\frac23\max(\lambda_m,\lambda_n)
 \qquad(m\ne n).
\]
Pairing the terms indexed by $(m,n)$ and $(n,m)$, their absolute sum is
bounded by a constant times
\[
 \sum_{m>n}
 \frac{1+\lambda_m}{\sqrt{mn\lambda_m\lambda_n}}
 (1+\lambda_m)^{-M},
\]
which is finite for $M\ge2$ because $\lambda_j=3^j$.  This bound is
independent of $N$.  Thus the left side of
\eqref{eq:divergent-current} tends to $+\infty$, although $f_N$ converges in
$L^2$.  This does not give a phase-retrieval counterexample.  It shows that
the current required by the $H^1$ proof cannot be obtained by an arbitrary
smooth $L^2$ approximation.

The other regularizations have the same defect.  A spatial convolution does
not preserve the equation with the same potential and does not preserve the
modulus identity.  Spectral cutoffs, resolvents, and the heat operators
\[
 \1_{[m,N]}(H),
 \qquad
 (1+\varepsilon(H-m))^{-1},
 \qquad
 e^{-\varepsilon(H-m)}
\]
produce finite-energy vectors, but in general
\[
 |u|=|v|\quad\not\Longrightarrow\quad
 |T(H)u|=|T(H)v|.
\]
Time convolution is a spectral multiplier and has the same problem.  Real
Schr\"odinger evolution supplies no spatial smoothing: already for $V=0$,
$e^{it\partial_x^2}f_0$ belongs to $H^s$ at one time if and only if
$f_0\in H^s$.

\subsection*{The exterior-product route}

One part of the argument remains valid at the bare $L^2$ level.  Define
\[
 F(t,x,y)=u(t,x)v(t,y)-v(t,x)u(t,y).
\]
Then $F\in C(\R;L^2(\R^2))$ and
\[
 iF_t+F_{xx}+F_{yy}=(V(x)+V(y))F
 \tag{13.3}\label{eq:masuda-exterior-equation}
\]
in distributions.  The lifted Hamiltonian is
\[
 H^{(2)}=H\otimes\Id+\Id\otimes H.
\]
It is self-adjoint and semibounded.  The algebraic tensor product of the
one-dimensional core with itself is a core for $H^{(2)}$, and
$V(x)+V(y)\in L^2_{\mathrm{loc}}(\R^2)$.  The lifted exceptional set is
\[
 (Q\times\R)\cup(\R\times Q),
\]
whose complement is connected.  Thus Masuda's theorem applies in two spatial
dimensions to solutions of \eqref{eq:masuda-exterior-equation}.

In the diagonal coordinates
\[
 r=\frac{x+y}{\sqrt2},
 \qquad
 s=\frac{x-y}{\sqrt2},
\]
the function $F$ is odd in $s$.  Its formal Dirichlet trace at $s=0$ is zero,
but its formal normal derivative is a nonzero constant multiple of
\[
 vu_x-uv_x.
 \tag{13.4}\label{eq:missing-masuda-wronskian}
\]
If one could define the two currents and prove
\[
 \overline u\,u_x=\overline v\,v_x,
 \qquad vu_x-uv_x=0,
\]
at the critical endpoint, the scalar trace argument would give zero
Dirichlet and Neumann traces.  The extension
\[
 G=\1_{\{s>0\}}F
\]
would then solve \eqref{eq:masuda-exterior-equation} and vanish on an open
half-space.  After extending this local weak solution by the unitary group for
$H^{(2)}$, Masuda's theorem would give $G=0$, hence $F=0$.  The rank-one
argument would finish the proof.

This route identifies the exact missing estimate.  One needs enough
real-time endpoint regularity, together with a valid bounded-antiderivative
continuity identity, to define and annihilate
\eqref{eq:missing-masuda-wronskian}.  For globally bounded-plus-square-integrable potentials, Strichartz and
local-smoothing estimates~\cite{KeelTao,VegaVisciglia} provide the required
endpoint identities.  They are
not consequences of Masuda's local assumptions and semiboundedness alone in
the present argument.

We therefore retain the $H^1$ hypothesis in Theorem~\ref{thm:phase-retrieval}.
No counterexample under all of Masuda's assumptions is known here, so the
corresponding $L^2$ phase-retrieval statement is not declared false.  It
remains open in this manuscript.  For $V\in L^\infty(\mathbb R)+L^2(\mathbb R)$, the separate
common-potential argument based on Strichartz estimates proves the $L^2$
conclusion from measurements on any nonempty open time interval.

Finally, Masuda's theorem concerns vanishing on an open subset of spacetime.
Compact support at one isolated time does not imply that a Schr\"odinger
solution is zero.  In the finite-energy proof, the density identity first
produces the open set $\Omega_0$ on which $v-\zeta u$ vanishes; Masuda's
theorem is applied only after that step.

\begin{thebibliography}{9}

\bibitem{DeiftTrubowitz}
P.~Deift and E.~Trubowitz,
\emph{Inverse scattering on the line},
Comm. Pure Appl. Math. \textbf{32} (1979), 121--251.

\bibitem{EgorovaKopylovaMarchenkoTeschl}
I.~Egorova, E.~Kopylova, V.~Marchenko, and G.~Teschl,
\emph{Dispersion estimates for one-dimensional Schr\"odinger and
Klein--Gordon equations revisited},
Russian Math. Surveys \textbf{71} (2016), 3--26.

\bibitem{TeschlBook}
G.~Teschl,
\emph{Mathematical Methods in Quantum Mechanics},
2nd ed., Graduate Studies in Mathematics 157,
American Mathematical Society, 2014.

\bibitem{KeelTao}
M.~Keel and T.~Tao,
\emph{Endpoint Strichartz estimates},
Amer. J. Math. \textbf{120} (1998), 955--980.

\bibitem{VegaVisciglia}
L.~Vega and N.~Visciglia,
\emph{On the local smoothing for the Schr\"odinger equation},
Proc. Amer. Math. Soc. \textbf{135} (2007), 119--128.

\bibitem{Masuda}
K.~Masuda,
\emph{A unique continuation theorem for solutions of the Schr\"odinger
equations},
Proc. Japan Acad. \textbf{43} (1967), 361--364.

\end{thebibliography}

\end{document}
